\documentclass[12pt, a4paper]{report}
%\documentclass[11pt, a4paper]{article}

%====================== PACKAGES ======================
\usepackage[french]{babel}

\frenchbsetup{StandardLists=true}
\usepackage{enumitem}
\usepackage{pifont}

\usepackage[utf8x]{inputenc}
%\usepackage[latin1]{inputenc}

%pour gérer les positionnement d'images
\usepackage{float}
\usepackage{amsmath}
\DeclareMathOperator{\dt}{dt}
\usepackage{graphicx}
%\usepackage{tabularx}
\usepackage[colorinlistoftodos]{todonotes}
\usepackage{url}

%pour les informations sur un document compilé en PDF et les liens externes / internes
\usepackage[pdfborder=0]{hyperref}
\hypersetup{
	colorlinks = true
	}

%pour la mise en page des tableaux
\usepackage{array}
\usepackage{tabularx}
\usepackage{multirow}
\usepackage{multicol}
\setlength{\columnsep}{50pt}

%pour utiliser \floatbarrier
%\usepackage{placeins}
%\usepackage{floatrow}

%espacement entre les lignes
\usepackage{setspace}

%modifier la mise en page de l'abstract
\usepackage{abstract}

%police et mise en page (marges) du document
\usepackage[T1]{fontenc}
\usepackage[top=2cm, bottom=2cm, left=2cm, right=2cm]{geometry}

%Pour les galerie d'images
\usepackage{subfig}

\usepackage{pdfpages}

\usepackage{tikz}
\usetikzlibrary{trees}
\usetikzlibrary{decorations.pathmorphing}
\usetikzlibrary{decorations.markings}
\usetikzlibrary{decorations.pathreplacing,calligraphy}
%\usetikzlibrary{decorations}
\usetikzlibrary{angles, quotes}
\usepackage{verbatim}

\usepackage{appendix}

\usepackage{comment}

\usepackage{xcolor}

%\PreviewEnvironment{tikzpicture}
%\setlength\PreviewBorder{0pt}%

%====================== INFORMATION ET REGLES ======================

%rajouter les numérotation pour les \paragraphe et \subparagraphe
\setcounter{secnumdepth}{4}
\setcounter{tocdepth}{4}

\hypersetup{							% Information sur le document
pdfauthor = {Stephan Runigo},			% Auteurs
pdftitle = {Documentation},			% Titre du document
pdfsubject = {Documentation},		% Sujet
pdfkeywords = {Document},	% Mots-clefs
pdfstartview={FitH}}	% ajuste la page à la largeur de l'écran
%pdfcreator = {MikTeX},% Logiciel qui a crée le document
%pdfproducer = {} % Société avec produit le logiciel

%======================== DEBUT DU DOCUMENT ========================
%
\begin{document}
%
%régler l'espacement entre les lignes
\newcommand{\HRule}{\rule{\linewidth}{0.5mm}}
%
% Titre, résumé, ... %
%
%
\begin{titlepage}
%
~\\[1cm]

\begin{center}
%\includegraphics[scale=0.5]{./presentation/chambreABulle}
\end{center}

\textsc{\Large }\\[0.5cm]

% Title \\[0.4cm]
\HRule

\begin{center}
{\huge \bfseries Spiritualisme et\\
matérialisme\\[0.4cm] }
\end{center}

\HRule \\[1.5cm]


\vfill

\hfill
\begin{minipage}{0.4\textwidth}
\begin{flushright} \large
%\emph{Auteur:}\\
%Stephan \textsc{Runigo}
Extraits de dictionnaires et d'encyclopédies
\end{flushright}
\end{minipage}

\vfill
{\sf \footnotesize
\begin{itemize}[leftmargin=1cm, label=\ding{32}, itemsep=1pt]
\item {\bf Objet : } Étudier les concepts liés au spiritualisme et au matérialisme.
\item {\bf Contenu : } Définitions et philosophie encyclopédique.
\item {\bf Public concerné : } Néophyte.
\end{itemize}
}

\vfill

% Bottom of the page
{\large \today}

\end{titlepage}

\newpage
%\begin{center}
\Large
Résumé
\normalsize
\end{center}
\vspace{3cm}
\begin{itemize}[leftmargin=1cm, label=\ding{32}, itemsep=21pt]
\item {\bf Objet : éclairer } .
\item {\bf Contenu : définitions} .
\item {\bf Public concerné : néophyte} .
\end{itemize}

\vspace{3cm}

résumé.

\vspace{3cm}

Ce document contient des définitions provenant de divers dictionnaires.




\thispagestyle{empty}

\begin{center}
\Large
%Introduction
Préambule
\normalsize
\end{center}
\vspace{3cm}

Ce document est une compilation d'articles provenant de quatre ouvrages : un dictionnaire encyclopédique de poche, {\it La pratique de la philosophie} destiné aux lycéens, une encyclopédie de la philosophie destinée aux néophytes, et le dictionnaire philosophique d'André Comte-Sponville.
les chapitres contiennent les articles des trois premiers ouvrages, les articles du dictionnaire de philosophie d'André Comte-Sponville sont reproduits en annexe.

\vspace{1.3cm}

Chaque chapitre contient les articles correspondant à une notion particulière. Ces notions ont été choisies en raison de leurs liens avec la question du hasard. Ces choix ont été guidés : 1. Par les renvois vers d'autres articles présent dans les ouvrages, 2. Mes propres choix, liés à mes connaissances, 3. La volonté d'obtenir une quantité raisonnable d'information.

\vspace{1.3cm}

Dans {\it La pratique de la philosophie}, l'article concernant la nécessité renvoit à des textes de Spinoza dont le choix reste subjectif à l'ouvrage. J'ai néanmoins reproduit en annexe ces textes ainsi que l'article concernant Spinoza. (les autres annexes sont d'autres renvois de cet ouvrage)

Les articles compilés dans ce document comportent donc les choix "discutables" réalisés dans les trois ouvrages utilisés. Il s'agit donc d'un document de "travail" destiné à apporter quelques points de vues philosophiques de manière relativement élémentaire.

\vspace{1.3cm}

Les trois premiers chapitres abordent les thèmes du hasard, de la nécessité puis du vitalisme. Les chapitres suivants élargissent le champ de vision philosophique en abordant les thèmes du déterminisme, de la contingence et de la providence.

\vspace{2.3cm}

\hfill Stephan Runigo

%%%%%%%%%%%%%%%%%%%%%%%%%%%%%%%%%%%%%%%%%%%%%%%%

%
\newpage
%
%

\thispagestyle{empty}
\begin{center}
\Large
%Introduction
Les ouvrages utilisés
\normalsize
\end{center}

\vspace{1.3cm}

Premières de couverture

\begin{center}
\includegraphics[scale=0.43]{./presentation/dictionnaire-1}
\hfill
\includegraphics[scale=0.43]{./presentation/pratique-1}
\hfill
\includegraphics[scale=0.43]{./presentation/encyclopedie-1}
\end{center}

\vspace{1.3cm}

Quatrièmes de couverture

\begin{center}
\includegraphics[scale=0.43]{./presentation/dictionnaire-2}
\hfill
\includegraphics[scale=0.43]{./presentation/pratique-2}
\hfill
\includegraphics[scale=0.43]{./presentation/encyclopedie-2}
\end{center}



%\newpage
%

%
% Table des matières
\tableofcontents
\thispagestyle{empty}
\setcounter{page}{0}
%
%espacement entre les lignes des tableaux
\renewcommand{\arraystretch}{1.5}
%
%====================== INCLUSION DES CHAPITRES ======================
%
~
\thispagestyle{empty}
%recommencer la numérotation des pages à "1"
\setcounter{page}{0}
\newpage
%
\chapter{La philosophie}
%

%%%%%%%%%%%%%%%%%%%%%
\section{De l’ « idée » à l’ « impression »}
%%%%%%%%%%%%%%%%%%%%%
%
%{\it }
La philosophie anglaise du {\footnotesize XVIII}$^\text{e}$ siècle, amoureuse du concret, du donné, méfiante à l’égard des abstractions et
des hypothèses, se veut tout à la fois ingénue et rigoureuse ; rigoureuse parce qu’ingénue, parce qu’attentive à
rester fidèle à l'expérience vécue, à l’immédiat.

Ceci est plus sensible encore chez Hume que chez
Berkeley. L’ingénuité berkeleyenne n’est pas exempte
d'artifice et de présupposés. Toute l’œuvre de Berkeley
est inspirée par des préoccupations apologétiques. Hume,
en revanche, est un pur philosophe qui entend se laisser
conduire jusqu’au bout de sa réflexion quelles qu’en
soient l’issue et les conséquences. Discutant sur le problème
de la liberté il note fermement qu’ « il n’est pas certain
qu’une opinion soit fausse de ce qu’elle est de dangereuse
conséquence ».

Le radicalisme philosophique de Hume se manifeste
d’abord en ce qu’il entend remonter à des données originaires. Le point de départ de la réflexion philosophique
se trouvera donc dans ces données de la conscience, au
sens le plus large que Locke et Berkeley nommaient les
{\it idées} et que Hume appelle les {\it perceptions}. Car tandis que
Locke et Berkeley entendent par idées tout ce qui constitue
le contenu de la conscience, Hume opère une distinction
importante. La toute première phrase du {\it Traité de la nature
humaine} dit en effet que « toutes les perceptions de
l'esprit humain se ramènent à deux genres distincts que
%{\it }
%18
j’appellerai impressions et idées ». Seules les impressions
sont originaires, les idées ne sont que « des copies de nos
impressions », des reflets atténués de nos sensations dans
le miroir de nos pensées. Sommairement présentée, cette
thèse semble être celle de l’empirisme le plus plat. La
doctrine de Hume s’identifierait à un sensualisme grossier
et se contenterait de reprendre la formule célèbre de
Locke : « Il n’y a rien dans l’entendement qui n’ait d’abord
été dans les sens. {\it Nihil est in intellectu quod non prius fuerit in sensu.} »

Mais si nous l’examinons avec quelque attention, nous
voyons que la distinction humienne des impressions et
des idées se révèle beaucoup plus intéressante et beaucoup
plus profonde.

Tout d’abord il ne faudrait pas imaginer que pour
Hume l'esprit soit purement passif, une « table rase »,
une pâte molle où s’inscriraient mécaniquement les stimuli
externes. Assurément un aveugle-né ne saurait, faute
d’impressions, se faire la moindre idée des couleurs. Mais
supposons « un homme familiarisé avec les couleurs de
tout genre sauf avec une nuance particulière de bleu que
le hasard ne lui a jamais fait rencontrer. Qu’on place
devant cet homme toutes les diverses nuances de cette
couleur, à l’exception de cette nuance particulière, dans
une gradation descendante de la plus foncée à la plus
claire ». Cet homme, assure Hume, « percevra un vide »,
pourra « suppléer à ce défaut par sa seule imagination »
et « se donner l’idée de cette nuance particulière que
cependant ses sens ne lui ont jamais fournie ». Exception
unique mais très significative. Ceci témoigne d’un élan
de l'imagination, d’un dynamisme de l'esprit humain,
d’une activité psychologique subjective qui, dans l’empirisme original du philosophe écossais, est fondamentale.

De plus l'impression ne s’oppose pas à l’idée comme
une sensation d’origine externe s’opposerait à un phénomène psychologique intérieur. À vrai dire Hume ne
s’interroge pas sur la source des impressions. Pour lui les
%19
impressions sont des données originaires au-delà
desquelles on ne saurait remonter. En ce sens Hume n’est
nullement l’adversaire de l’innéisme, il le remarque
expressément : « Si l’on entend par inné ce qui est primitif,
ce qui n’est copié d'aucune impression antérieure, alors
nous pouvons affirmer que toutes nos impressions sont
innées et que nos idées ne le sont pas. »

Il faut ajouter qu’à côté des « impressions de sensation »
nous avons des  « impressions de réflexion » qui comprennent
les émotions, les passions, tous les phénomènes de
désir et de volonté. Sans doute « la plupart » des impressions de réflexion proviennent indirectement des impressions de sensation par l'intermédiaire des idées. Par
exemple l’idée de plaisir, l’idée de douleur provoquent dans
l’âme de nouvelles impressions de désir ou d’aversion.
Mais il reste que la distinction impression-idée est une
distinction purement psychologique. Elle s’opère sur des
critères internes. L’impression se définit par ce que Hume
appelle sa {\it vividness}, sa {\it liveliness}, c’est-à-dire
moins peut-être par son intensité que par sa présence vivante, par son
caractère actuel. L’impression ne désigne donc pas spécifiquement
la sensation d’origine externe mais, comme dit
Hume, c’est « un mot tout nouveau pour désigner nos
perceptions les plus vivantes quand nous entendons,
voyons, touchons, aimons, haïssons et voulons ».

%{\it }
Les philosophes ne raisonnent que trop souvent sur des
« idées indistinctes et obscures ». Ils croient penser et ils
ne font que parler. Hume les invite alors à rechercher les
« impressions » authentiques. Il ne s’agit pas de tout
ramener au sensoriel. Il s'agit plutôt d'une critique du
langage, d’une invitation à retrouver la pensée vivante,
actuelle : « Quand nous soupçonnons qu’un terme philosophique
est employé sans aucun sens ni aucune idée
correspondante, nous n’avons qu’à rechercher {\it de quelle
impression dérive cette idée supposée}. » Aller de l’idée à
l'impression c’est demander aux philosophes « ce qu’ils
entendent par substance et inhérence » ; c’est leur demander :
%20
« Qu’existe-t-il effectivement ? » Il s’agit de « porter
les idées à la lumière » pour écarter les discussions stériles
en démasquant les pseudo-idées. L'appel à l’impression
signifie chez Hume, dit très bien Laporte, « non le préjugé
du sensualisme mais la haine du verbalisme ». Les impressions
de Hume ce sont les idées claires de Descartes, celles
qui sont « présentes et manifestés à un esprit attentif »,
ce sont les « données immédiates » que Bergson voudra
retrouver. Et il est très remarquable que Husserl invoquera
Hume, au même titre que Descartes, comme précurseurs
de sa méthode phénoménologique dônt le but sera précisément d’écarter les hypothèses et les théories pour
revenir « aux choses elles-mêmes », pour tenter d’élucider,
de dévoiler ({\it enthüllen}) des significations. Ce que nous
demande Husserl, commente Gaston Berger, « c’est de
revenir aux données originaires, à ce qui est supposé par
toutes les opérations symboliques et toutes les représentations mais que le jeu même de ces opérations symboliques
risque toujours de nous faire perdre de vue. Il faut aller
des concepts vides par lesquels une idée est seulement
visée à l’intuition directe et concrète de l’idée, tout comme
Hume nous apprend à revenir des idées aux impressions ».
Retrouver les impressions, c’est tout simplement retrouver l’immédiat.
Comme le dit Hume dans ses {\it Dialogues
sur la religion naturelle} : « Est-il nécessaire d’établir ce
que chacun éprouve au-dedans de soi ? Il est seulement
nécessaire de nous le faire éprouver, s’il se peut, d’une
façon plus intime et plus sensible. » Telle est la tâche difficile de la philosophie, car cette ingénuité qui nous est
demandée est, en fait, une conquête ardue. L’immédiat
n’est pas donné immédiatement ; il faut le reconquérir
en deçà des préjugés et des théories, en deçà des formules
vides qui s’insinuent et s'imposent, créant des habitudes
de langage qui se transforment en habitudes de pensée ou
plutôt de non-pensée. Le contact avec l'immédiat ne saurait
être un point de départ ; il est le fruit et la récompense
d’une critique sévère.
%21
%%%%%%%%%%%%%%%%%%%%%%%%%%%%%%%%%%%%%%%%%%%%%%%%%%%%%%%%%%%%%%%%%%%%%%%%

%

%%%%%%%%%%%%%%%%%%%%%
\section{Critique de quelques notions énigmatiques : causalité, probabilité, substance}
%%%%%%%%%%%%%%%%%%%%%
%{\it }
L’analyse de la relation de causalité par Hume est
justement célèbre. Cette relation où l’on voit communément le principe fondamental de la raison humaine est
pour un philosophe empiriste assez extraordinaire. Tandis
qu’une relation d’idées comme la ressemblance porte sur
plusieurs faits qui tous sont donnés dans l’expérience
— et que nous pouvons alors comparer sans difficulté —,
la relation causale porte sur deux termes dont l’un est
absent : « Pas de fumée sans feu. » Je constate la présence
de la fumée et j’en infère au nom du principe de causalité
qu’un feu a été allumé ; ou bien je constate qu’on allume
un feu et je prévois qu’il y aura de la fumée. Au nom de
la relation causale, je « dépasse le témoignage de mes
sens », je prévois ce qui va se produire, j’infère quelque
chose dont je n’ai pas l’expérience actuelle. De quel droit
puis-je tirer d’un fait donné une conséquence qui le déborde,
de quel droit mon esprit fait-il ce « pas en avant » ? On
répond en invoquant la « connexion nécessaire » par laquelle
tel fait donné produira infailliblement tel autre fait. On
parle de « pouvoir », d’ « efficace », de « qualité productrice ». Il y aurait entre le fait présent et ce qu’on en
infère un lien nécessaire tel que le fait étant donné, l’autre
fait ne peut pas ne pas se produire. Mais d’où vient cette
idée de connexion nécessaire qui est, en effet, le nerf du
principe de causalité ? Il faut ici, selon la méthode de
Hume, découvrir l’impression originaire dont cette idée
dérive. Et pour trouver cette impression il faut se livrer
à ce que Hume appelle une enquête, il faut chercher cette
impression « à toutes les sources d’où elle peut découler »
comme on organise, dit Hume, une battue dans une forêt
pour trouver un objet précieux qui a été perdu.
%22
%{\it }
Demandons-nous d’abord si l’idée de cause peut dériver
d’une impression de sensation.

Tout le monde peut constater que la fumée apparaît
quand le feu est allumé, que l’ébullition suit l’échauffement
de l’eau, que lorsqu'une boule de billard en heurte une
autre, cette seconde boule se meut à son tour. Mais cette
constatation banale est celle d’une simple conjonction,
d’une succession entre deux phénomènes. Nous voyons
bien que l’échauffement de l’eau a précédé l’ébullition,
que le mouvement de la première boule de billard a précédé
le mouvement de la seconde. Mais nous ne voyons pas
entre les deux événements de connéxion nécessaire. Nous
constatons un « et puis », nous ne constatons pas un
« parce que ». Nous avons beau examiner, retourner en
tous sens le phénomène que nous appelons « cause » nous
ne pouvons découvrir en lui l'« efficace » qui produirait
l'effet. Rien dans le mouvement de la première bille de
billard ne nous donne la moindre indication sur le mouvement
qui va lui succéder. Nul objet dit Hume, « ne nous
montre jamais par les qualités qui paraissent aux sens,
soit les causes qui le produisent, soit les effets qui en
naissent ». Aucun effet n’est assignable a {\it priori}. Il est
vrai que nous pouvons répéter l'expérience, que nous
avons pu constater des centaines de fois, que l’échauffement de l’eau provoque son ébullition. La conjonction
entre les deux événements apparaît dans notre expérience
passée comme une conjonction constante. Mais le pourquoi
de cette liaison ne nous est toujours pas donné. La répétition
à l'infini d’une énigme n’est pas une solution de cette
énigme : « La simple répétition d’une impression passée,
même à l'infini, n’engendrera jamais une nouvelle idée
originale comme celle de connexion nécessaire ; le nombre
d’impressions n’a, dans ce cas, pas plus d'effet que si nous
nous en tenions à une seule. » L'idée de connexion nécessaire ne peut donc provenir d’une impression de sensation.

Tournons-nous à présent vers les impressions de réflexion,
c’est-à-dire vers le sentiment que nous avons des opérations de notre esprit.
%23
%{\it }
N’avons-nous pas la conscience
immédiate d’être la cause des mouvements de notre
corps propre ? (C’est ce que Maine de Biran affirmera
au début du {\footnotesize XIX}$^\text{e}$ siècle.) N’avons-nous pas également le
sentiment de notre pouvoir efficace sur nos volitions, sur
le cours de nos idées ? (C’est ce que soutiendra W. James
à la fin du {\footnotesize XIX}$^\text{e}$ siècle.) Mais Hume va répondre négativement 
à ces deux questions. Nous avons bien l’impression,
par exemple, d’une succession, d’une conjonction constante
entre notre intention d’accomplir un mouvement, et ce
mouvement lui-même. Nous constatons que nous voulons
et puis que le mouvement s’effectue. Mais ici encore les
deux événements nous sont donnés dans l’expérience
sans que nous saisissions en nous un véritable « pouvoir
efficace ». Nous constatons que la volonté paraît avoir un
effet sur la langue ou sur les doigts, non sur le cœur et sur
le foie : « Cette question ne nous embarrasserait pas si
nous avions conscience d’un pouvoir dans le premier
cas. » L’homme qui vient d’être frappé de paralysie est
tout surpris de ne pouvoir remuer une jambe. Il ressent
son vouloir comme nous ressentons le nôtre. Mais son
vouloir n’est suivi d’aucun effet. Notre vouloir est bien
suivi d’un effet, mais nous ne comprenons pas plus notre
pouvoir que le paralytique ne comprend son impuissance.
D'ailleurs « l'anatomie » nous enseigne que le mouvement
ne se produit que par le jeu des os et des muscles, par
l’action des « esprits animaux » que nous n’avons aucunement 
conscience de déclencher. Si nous réfléchissons
honnêtement, nous saisissons que le lien entre notre désir
et les mouvements de notre corps reste aussi mystérieux
que si nous avions «le pouvoir d’écarter les montagnes ou
de contrôler les astres dans leurs orbites ».

La maîtrise de l'esprit sur lui-même n’est pas plus
claire que sa maîtrise sur le corps. Nous constatons sans
le comprendre que nos pensées nous obéissent mieux que
nos sentiments, que « nous sommes plus maîtres de nos
pensées le matin que le soir, à jeun qu’après un repas
%24
%{\it }
copieux ». « Où est alors le pouvoir dont nous prétendons
être conscients ? » Dans tous les cas nous avons seulement
des impressions de contiguité et de succession entre les
événements. Et cependant l’idée de causalité est tout autre
chose : « Un objet peut être contigu et antérieur à un autre
sans qu’on le considère comme sa cause. » Mais cette idée de
causalité, de connexion nécessaire n’est nullement éclaircie.

Tout au long de cette enquête Hume a suivi de très
près la critique de Malebranche. Celui-ci concluait, comme
on sait, que l’efficace n’est nullement dans le phénomène
abusivement appelé cause et qui n’est qu’une occasion.
Malebranche refuse aux créatures, même les plus élevées,
tout pouvoir causal (« un ange ne saurait remuer un
fétu de paille ») et place la causalité en Dieu seul. Cette
hypothèse métaphysique est pour Hume « trop étrange
pour jamais apporter avec elle la conviction à aucun
homme suffisamment informé de la faiblesse de la raison
humaine ». Les grands systèmes métaphysiques nous
proposent un voyage « au pays des fées »; ce sont des
fictions qui débordent le domaine de l’expérience, des
impressions réellement données : « Notre ligne est trop
courte pour sonder l’immensité de pareils abîmes. »
L'analyse de Hume va s’orienter dans un domaine tout
différent. Elle va se placer à un point de vue nouveau qui
est le point de vue psychologique. Hume entend bien ne
pas sortir du domaine des faits et le seul fait incontestable
est ici un fait psychologique. Bien que l'« enquête » philosophique 
ne nous ait livré aucune impression originale
de causalité, c’est un fait que nous croyons tous à la cause,
à la connexion nécessaire, au pouvoir efficace. C’est donc
notre croyance qu’il faut expliquer. Nous n’avons l’expérience 
d’aucune force efficace présente dans un phénomène,
mais pourquoi croyons-nous que ce phénomène sera nécessairement 
suivi d’un autre ? Il ne s’agit pas de fonder
rationnellement le principe de causalité, il s’agit de rendre
compte de notre croyance. Pour résoudre le problème,
Hume le réduit à sa dimension psychologique. Il ne s’agit
%25
%{\it }
pas tant de chercher le fondement légitime du principe
de causalité que d’élucider son origine psychologique.

Or, la constante conjonction des objets, nous l’avons
dit, n’a aucune influence sur les objets eux-mêmes. Le
millième cas de conjonction est tout aussi énigmatique
que le premier. Seulement cette conjonction constante a
de l'influence sur notre esprit. L’habitude de voir deux
objets liés ensemble produit en nous une tendance très
forte à attendre le deuxième objet si le premier se présente
à nous une fois de plus : « Après une répétition de cas
semblables, l’esprit est porté par habitude lorsqu’apparaît
un des deux événements à attendre son concomitant
ordinaire et à croire qu’il existera. » C’est la transition,
le glissement aisé de l’imagination d’un objet à son concomitant 
habituel qui fournit l’unique impression d’où
dérive l’idée de connexion nécessaire. Dès lors celle-ci
n’a aucune portée ontologique, elle n’a qu’un sens 
psychologique : « Elle n’est qu’une impression intérieure de
l'esprit, une détermination à porter nos pensées d’un objet
à l’autre. » La causalité n’est donc pas un principe qui
régit les choses, elle n’est qu’un « principe de la nature
humaine ». Hume le déclare expressément : « La nécessité
est quelque chose qui existe dans l’esprit, non dans les
objets. » Hume a eu le sentiment très vif de faire œuvre
révolutionnaire par son analyse de la causalité. Désormais,
dit-il, « la connexion nécessaire dépend de l’inférence au
lieu que ce soit l’inférence qui dépende de la connexion
nécessaire ». Et il a bien saisi que ce renversement des
perspectives pouvait paraître paradoxal et scandaleux :
« Quoi ! l'efficacité des causes se trouve dans la détermination 
de l’esprit ! Comme si les causes n’opéraient pas en
toute indépendance de l'esprit et ne continueraient pas
d’opérer même s’il n’y avait aucun esprit pour les contempler 
et raisonner à leur sujet. La pensée peut bien dépendre
des causes pour son opération mais non les causes de la
pensée ! C’est renverser l’ordre naturel ! »

Ce renversement paradoxal n’en est pas moins la clef
%26
%{\it }
de toutes les réductions psychologiques opérées par Hume
à propos de toutes les idées qu’il soumet à son enquête.
La notion de probabilité subit le même déplacement
psychologique, qu’il s’agisse d’une probabilité confuse
lorsque notre inférence manque de conviction, parce que
l'expérience passée qui la détermine est trop éloignée ou
n’a pas été assez souvent répétée, ou qu’il s’agisse d’une
probabilité « philosophique » c’est-à-dire mathématique,
susceptible de se calculer : Par exemple en jetant un dé
à six faces, nous disons qu’il y a une chance sur six que
telle face (celle qui porte le n° 3 par exemple) sorte. Il
n’y a certes « rien de semblable à la Chance dans le monde »
et la probabilité est — comme la causalité — une impression ou plus exactement une impulsion de notre esprit.
Seulement l'élan de l’imagination qui, dans l’expérience
de la causalité est dirigé par l'habitude dans une seule
direction, se divise ici en autant de tendances qu’il y a
de faces sur le dé. Si quatre faces portent le même chiffre,
les deux autres étant différentes, nous disons que le chiffre
des premières a quatre chances sur six de sortir parce que
« l'impulsion moindre détruit l’impulsion plus forte dans
la mesure de sa propre force ».

Hume applique la même méthode psychologique dans sa
critique de la notion de substance. Il retrouve ainsi à propos
de la substance matérielle le point de vue de Berkeley. Le
sens commun croit à l’existence de «substances », à la permanence 
des choses bien que nous apercevions constamment
de nouveaux objets. Quand je sors de mon bureau, je ne vois
plus la table, la bibliothèque, le tapis, mais je ne doute
pas de les retrouver tout à l’heure. Le sens commun admet
que « les perceptions discontinues sont reliées par une
existence réelle dont nous n’avons pas conscience ». Cette
bibliothèque vitrée en noyer est une chose qui existe,
même si je ne la vois pas ! Une telle affirmation, cependant,
n’est pour le philosophe rigoureux qu’une hypothèse
gratuite. Mon expérience me livre des phénomènes discontinus, rien de plus, et on ne saurait fonder en raison
%27
%{\it }
l'affirmation des « substances ». Hume substitue à cette
affirmation ontologique une simple explication psychologique 
de notre croyance aux substances. En fait notre
imagination fond les unes dans les autres les apparences
semblables et successives. L’habitude de retrouver constamment 
les mêmes phénomènes à leur place engendre une
attente irrésistible. L’affirmation de l'existence de la
bibliothèque n’est que l'expression de la certitude que
j'ai de la retrouver en pénétrant dans mon cabinet de
travail. Mon imagination glisse si aisément le long de la
succession des apparences qu’elle métamorphose la simple
succession en identité véritable.

Mais Hume va beaucoup plus loin que Berkeley, car il
étend à la notion de substance spirituelle la critique de la
notion de substance matérielle. Le sens commun affirme
l'existence de substances spirituelles, il m’attribue par
exemple un moi un et identique, ce moi n’étant nullement
une impression mais ce à quoi nos diverses impressions ou
idées sont censées se rapporter. Mais je n’ai pas réellement
l'expérience d’un tel sujet permanent. J’assiste seulement à
l’incessant défilé de mes états de conscience. Souffrances et
joies, passions diverses, se succèdent sans jamais pouvoir
exister toutes ensemble en même temps. Le support de tous
ces phénomènes, cette « substance » qu’on appelle l'âme ou
le moi n’est à la lettre qu’une invention de notre imagination.
C’est ici encore l’imagination qui prolonge l’élan que la
mémoire lui communique et « comme une galère mise en
mouvement par les rames court sur son erre sans nouvelle
impulsion ». Habile à masquer la discontinuité de tous les phénomènes, elle nous persuade, en glissant aisément d’un état
à un autre, de l'identité et de la substantialité de notre être.

%%%%%%%%%%%%%%%%%%%%%%%%%%%%%%%%%%%%%%%%%%%%%%%%%%%%%%%%%%%%%%%%%%%%%%%%%%%

%

%%%%%%%%%%%%%%%%%%%%%
\section{Scepticisme et dogmatisme}
%%%%%%%%%%%%%%%%%%%%%
%{\it }
Ces analyses nous permettent, semble-t-il, de prendre
parti dans la controverse toujours ouverte du « scepticisme » 
de Hume. Laporte penchaït pour la thèse du
scepticisme absolu que rejetait jadis Compayré, que
conteste aujourd’hui M. Leroy qui verrait plutôt dans la
philosophie de Hume un « probabilisme ».

Hume assigne des limites à notre savoir, tous les historiens 
l’accordent, mais où ces limites se placent-elles
exactement ?

Hume distingue deux types de connaissance : les relations 
d’idées et les faits. Les relations d’idées (une proposition 
telle que trois fois cinq est égal à la moitié de trente
exprime une relation entre ces nombres) présentent cette
particularité que l’on peut les découvrir « par la seule
opération de la pensée sans dépendre de rien de ce qui
existe dans l’univers ». Les vérités nécessaires que nous
pouvons poser dans ce domaine sont des vérités purement
formelles, rigoureuses mais vides. Les conclusions auxquelles 
nous aboutissons ne sont rigoureusement nécessaires 
que parce qu’elles n’expriment rien d’autre que
l'accord de l'esprit avec ses propres définitions. Tel est
le propre de la pensée logique, mathématique, et aussi
juridique (une proposition telle que celle-ci : « Où il n’y
a pas de propriété il ne peut pas y avoir d’injustice »,
est une proposition nécessaire si nous avons préalablement
défini l'injustice comme la violation de la propriété).

Les propositions relatives aux faits sont bien différentes.
Si les données immédiates, les « impressions » originaires
s'imposent comme telles, on n’en peut dire autant des
conclusions que nous posons en raisonnant sur les faits.
Le principe de causalité — nerf du raisonnement expérimental — 
n’a rien à voir, nous l’avons montré, avec la
relation formelle de principe à conséquence. Il peut donc
sembler que Hume conteste sa valeur et se montre sceptique 
à l'égard de toutes nos inductions expérimentales.
Hume inaugurerait ainsi une tradition (promise en Angleterre 
à une grande fortune) de logiciens sceptiques selon
laquelle on peut concevoir une {\it pensée} vraie (la pensée
logique, purement formelle) mais non une {\it connaissance}
vraie (toute affirmation sur les faits devant se limiter à
%29
%{\it }
l’expérience immédiate). Cependant la plupart des historiens 
voient en Hume un partisan résolu de la méthode
expérimentale. Selon Norman Kemp Smith, Hume serait
sceptique seulement pour tout ce qui est artificiel, trop
subtil, trop compliqué. Pour Greig, il s’agit d’un « scepticisme
à limites bien définies. Il s’étend uniquement à
la théorie de la connaissance, à la métaphysique, il ne
touche pas à la morale, la politique, la méthode expérimentale, 
la vie commune ». Et on peut rappeler que déjà
Hegel lorsqu’il opposait au scepticisme antique (qui est
au fond un platonisme, qui ne doute des faits concrets
que pour nous amener à la contemplation des idées) le
scepticisme moderne (qui, à l’opposé, doute des idées
métaphysiques mais fait confiance à la raison expérimentale), 
proposait Hume comme exemple typique de ce
scepticisme moderne, qui serait donc seulement ce que
nous appelons aujourd’hui un positivisme. Auguste Comte
n’a-t-il pas signalé Hume comme un de ses précurseurs ?

Et Hume lui-même s’est fait l’apôtre de la science
expérimentale. Il proclame son admiration pour Newton
dont l’œuvre géniale « passera, triomphante, jusqu’à la
postérité la plus reculée ». Loin de suspecter, semble-t-il,
la méthode expérimentale il veut la généraliser en l’appliquant 
à ce que nous appelons aujourd’hui les « sciences
humaines ». Le {\it Traité de la nature humaine} est expressément
présenté comme un {\it Essai pour introduire la méthode expérimentale 
de raisonnement dans les sujets moraux}.

Seulement le projet de Hume en s’accomplissant aboutit
paradoxalement à mettre en question les principes mêmes
qui l’ont inspiré. L’étude de la « nature humaine » et plus
particulièrement des mécanismes de la croyance se présente, 
si l’on veut, comme une psychologie positive mais
elle débouche sur tout autre chose qu’une confirmation
et qu’une extension de la méthode expérimentale. En
fait, l'étude psychologique des origines de notre croyance
en la causalité ruine la valeur de ce principe de causalité,
fondement de la raison expérimentale elle-même.

%30
%{\it }
L'analyse de la causalité n’est nullement équivalente
chez Hume, nous l’avons montré, à une théorie du fondement de 
l'induction. Le principe de causalité se réduit au
fait psychologique de l’attente d’un phénomène, attente
produite par une longue habitude. Le principe de causalité
ne peut donc pas fonder en raison nos habitudes inductives, puisqu'il 
n’est pas lui-même autre chose que l’expression. de ces habitudes. 
Expliquer, par exemple, notre
attente de l’ébullition de l’eau quand nous la chauffons,
par l’habitude ce n’est pas justifier notre attente. Rien
ne prouve que l’avenir ressemblera au passé, que le soleil
se lèvera demain, que l’eau chauffée demain se mettra à
bouillir. Il n'y a pas d’impossibilité {\it logique} que les arbres
se mettent à « fleurir en décembre et à dépérir en mai ».
Le contraire d’un fait est toujours possible car il n’implique pas contradiction. 
Le prince indien qui refusait de
croire les premières relations sur le gel ne raisonnait pas
autrement que les maîtres de la connaissance expérimentale. 
N'ayant « jamais vu d’eau en Moscovie pendant
l’hiver », il induisait, comme tout le monde, d’après ses
habitudes. Il ne pouvait croire en l’existence de la glace
puisqu'il avait « toujours vu l’eau fluide sous son propre
climat ». Tout se passe en effet comme si « n’importe
quoi pouvait produire n’importe quoi », et Hume dans son
{\it Histoire d'Angleterre} s’applique à montrer que des événements 
dérisoires ont souvent pour effet des conséquences
capitales.

La découverte de la signification purement psychologique 
de la causalité n’est pas présentée par Hume comme
une victoire de la raison expérimentale, mais comme
une conclusion qui ruine la valeur de notre savoir : « Quelle
déconvenue éprouvons-nous nécessairement quand nous
apprenons que cette connexion, ce lien, cette énergie se
trouvent seulement en nous; que ce n’est rien qu’une détermination 
de l'esprit acquise par accoutumance. » Une
telle découverte « coupe court à jamais à l'espoir d’obtenir
satisfaction », elle futilise aussi bien les inductions de
%31
%{\it }
la science expérimentale que celles de la vie quotidienne.
« L’entendement se détruit complètement lui-même et
ne laisse plus le moindre degré d’évidence à aucune proposition 
de la philosophie ou de la vie courante. »

Le scepticisme philosophique de Hume se situe, on le
voit, sur un tout autre plan que le doute provisoire du
savant. Il arrive que le savant s'interroge longuement
sur la cause d’un phénomène qu’il ne sait pas rattacher
à ce qu’il appelle les lois de la nature. Mais le savant ne
doute pas un seul instant de la validité de ces lois, ou
tout au moins de l’existence de lois naturelles même s’il
les ignore. Quand une cause « ne produit pas son effet
habituel » le savant pense que « d’autres causes cachées
sont intervenues ». Le doute du savant n'intervient que
sur fond de dogmatisme, de confiance dans l’ordre de la
nature. Ce qui étonne Hume ce n’est pas tel phénomène
isolé, c’est l’ensemble des phénomènes, leur ordre habituel
mais incompréhensible dans la nature tout entière. Le
doute du savant porte, comme le doute du vulgaire, sur
{\it l’insolite}. Le doute du philosophe porte tout au contraire
sur le cours {\it habituel} du monde parce que le philosophe
a su reconnaître dans l’habitude un simple fait psychologique, qu’il a su se délivrer de cette magie de Phabitude qui chez la plupart des hommes endort l’étonnement.

Le scepticisme philosophique de Hume est donc un
scepticisme absolu qui n’épargne même pas les raisonnements de la vie quotidienne, et c’est cela qui explique
le ton tragique et l’accent de désespoir des pages qui terminent la quatrième partie du {\it Traité de la nature humaine}
(ton de désespoir qui surprend évidemment les historiens
partisans de la thèse d’un scepticisme modéré chez Hume).
Hume se juge alors « exclu de tout commerce humain,
complètement abandonné et sans consolation » : « Quand
je tourne mes regards en moi-même, je ne trouve que
doute et ignorance », « je suis effrayé et confondu de cette
solitude désespérée où je me trouve placé dans ma philosophie ».

%32
%{\it }
Cependant ce scepticisme demeure « philosophique »,
c’est-à-dire purement spéculatif. Il est sans prises et sans
influence sur notre pratique quotidienne. A l’opposé des
maîtres stoïciens dont il fut le fervent lecteur, Hume
juge qu’il n’est ni souhaitable, ni possible de vivre comme
on pense. Le scepticisme c’est l’attitude purement intellectuelle, 
purement réflexive, de quelques moments de
spéculation dans le cabinet de travail du philosophe.
Mais dès que Hume, étourdi par ses doutes, quitte son
cabinet, se détend, rejoint ses amis au club et « joue au
tric-trac », la « nature» le « guérit » de sa « mélancolie
philosophique ». Après quelques heures de distraction,
les spéculations philosophiques: de Hume lui paraissent
«si froides, si forcées, si ridicules que je ne pourrais trouver
le cœur d’y pénétrer tant soit peu ». Le philosophe a opéré
une critique radicale de la connaissance ({\it Knowledge}).
Mais l’homme demeure rivé à ses croyances ({\it belief}).
Le mécanisme psychologique, qui nous porte à croire à la
causalité en transportant à son concomitant habituel
l'impression vivace que nous donne le fait actuellement
donné, continue à jouer dans l’imagination du philosophe
lui-même. Ce mécanisme « conservera son influence aussi
longtemps que la nature humaine demeurera la même ».
Nous n’avons pas à « craindre que notre philosophie sape
jamais les raisonnements de la vie courante » parce que,
précisément, la croyance est un acte « de la partie sentante
plutôt que de la partie pensante » de notre être. La croyance
est au fond un fait biologique. C’est « une espèce d’instinct
naturel » que le raisonnement philosophique ne peut pas
empêcher. « La nature, par une nécessité absolue et incompréhensible 
nous a déterminés à juger aussi bien qu’à
respirer et à sentir. » Ne constatons-nous pas que les gens
les plus incultes, que les animaux eux-mêmes savent en
quelque sorte « induire », tirer les leçons de l’expérience ?
Scepticisme et dogmatisme ne se distribuent pas chez
Hume selon les régions de la connaissance (comme s’il
y avait par exemple un scepticisme métaphysique et un
%33
%{\it }
dogmatisme scientifique), mais plutôt selon les niveaux
de la réflexion. Le scepticisme réflexif (qui est absolu)
coexiste en quelque sorte avec un dogmatisme instinctif,
Nous sommes ici tout près de Pascal qui notait que « la
nature confond les pyrrhoniens » tandis que « la raison
confond les dogmatiques ». Pour Hume aussi « la nature
maintiendra toujours ses droits et prévaudra à la fin
sur tous les raisonnements abstraits ». Tout se passe en
quelque sorte comme si la nature avait préféré confier à
l'instinct une fonction aussi importante que celle qui
consiste à prévoir les « effets » à partir des « causes ». Le
scepticisme de Hume aboutit donc par un nouveau renversement 
à une sorte de confiance irrationnelle envers la
nature, cette divinité du {\footnotesize XVIII}$^\text{e}$ siècle, 
véritable « Providence laïcisée ». Il y aurait comme une « harmonie 
préétablie entre le cours de la nature et la succession de nos
idées » bien que cette harmonie reste incompréhensible et
que les forces qui gouvernent la nature nous soient « totalement 
inconnues ». Le dernier mot de Hume serait-il
affirmation d’un finalisme mystérieux ? « Ceux, dit-il,
qui se réjouissent à contempler les causes finales ont ici
ample occasion de manifester leur émerveillement. »
Hume rejoindrait-il le « sentimentalisme » des philosophes
écossais, ses compatriotes auxquels on l’oppose d’ordimaire ? 
Brunschvicg n’hésitait pas à l’affirmer : « Hume a
respiré dans le {\footnotesize XVIII}$^\text{e}$ siècle l’atmosphère de l’optimisme.
L’argumentation sceptique des considérants ne lui interdit
pas les conclusions par lesquelles invoquant la bonté de
la nature en faveur de nos croyances spontanées, il rejoindra
ou précédera le gros de l’armée écossaise ! »

Pourtant cet optimisme finaliste ne se substitue pas au
scepticisme. Paradoxalement il s’y ajoute fournissant une
sorte de philosophie de rechange à l’usage de la pratique
et des moments où il convient de vivre plus que de penser.
Mais le scepticisme réapparaît sitôt que le philosophe se
livre de nouveau au démon de la réflexion solitaire et
désintéressée. Hume ne se fait pour ainsi dire aucune
%34
%{\it }
illusion sur la valeur philosophique de ce dogmatisme
instinctif. Il conclut curieusement qu’au philosophe sceptique 
« négligence et inattention peuvent seules apporter quelque 
remède et c’est pour cette raison que je me
repose entièrement sur elles ».

%%%%%%%%%%%%%%%%%%%%%%%%%%%%%%%%%%%%%%%%%%%%%%%%%%%%%%%%%%%%%%%%%%%%%%%%%%%

%

%%%%%%%%%%%%%%%%%%%%%
\section{Esthétique, morale, politique}
%%%%%%%%%%%%%%%%%%%%%
%34
%{\it }
Nous abordons ici le domaine des jugements de valeur.
Hume leur applique la même méthode de réflexion critique
qu'aux jugements de réalité portant sur les {\it matters of fact}.
Ici encore toute affirmation à prétention ontologique se
trouve réduite à une simple « croyance », à une réaction
subjective de la « nature humaine ».

Par exemple ce que nous appelons la beauté n’est pas
une qualité inhérente aux choses. Elle ne réside pas dans
des relations objectives. Autrement Euclide n’eût pas
manqué d’en faire une des propriétés du cercle ! En fait,
la beauté n’est que dans l’âme qui la contemple. Elle
est une réaction émotionnelle de Îa nature humaine (tout
comme la connexion nécessaire n’était pas dans les choses
mais dans l’esprit).

Mais subjectiviser les valeurs esthétiques n'est-ce pas
les futiliser, nous vouer au scepticisme ? L’un verra de la
laideur dans ce qui semble beauté à un autre et si toutes les
appréciations sont subjectives, « rechercher la vraie beauté
ou la vraie laideur est une enquête aussi vaine que de
prétendre établir le vrai doux ou le vrai amer ».

Seulement le scepticisme auquel nous aboutissons par là
est purement théorique. En fait les hommes s’accordent
assez bien sur les valeurs esthétiques. Nous reconnaissons
tous, dit Hume à plusieurs reprises, qu’ « il n’y a pas
autant de génie chez Ogilby que chez Milton » (Ogilby,
1600-1676, est un plat traducteur de Virgile). Une admiration durable 
s’attache aux grandes œuvres. Tout platonisme de l’idée de beauté est rejeté.
%35
%{\it }
Mais la validité pratique de nos jugements de goût repose sur l’uniformité
de la nature humaine. Le subjectivisme de Hume ne fait
donc pas obstacle à ce que Laporte appelait le « dogmatisme
du sentiment ». Cependant l’uniformité des sentiments
humains n’est pas absolue. On parle d’un bon goût et
d un mauvais goût. Il faudrait donc dit Hume dans sa
{\it Dissertation sur la règle du goût} « trouver une règle par
laquelle il serait possible de réconcilier les sentiments
divers des hommes ou du moins d’obtenir une décision
qui confirme un sentiment ou en condamne un autre »
Finalement, en matière de goût esthétique, on ne se ficra
qu’à des juges compétents dont le sens du beau est particulièrement 
délicat, affiné par une longue expérience
Mais c’est là substituer au {\it consensus omnium} le {\it consensus
optimi cujusque}. Sans doute le jugement de ces {\it happy few}
peut être empiriquement appuyé par la confiance que
l’ensemble des hommes leur accorde. Mais Hume lorsqu’il
réfléchit en philosophe trouve cette justification fragile
En fait, parler d’un bon et d’un mauvais goût, c’est corriger
le subjectivisme par un finalisme discret qui reste une
hypothèse métaphysique. Hume oscille entre un optimisme
purement empirique et un scepticisme réfléchi. Le problème 
du juge compétent dans le cadre du subjectivisme
de la {\it Dissertation sur la règle du goût} demeure une « question
embarrassante qui semble nous replonger dans la même
incertitude dont le but de cette dissertation était de nous
délivrer ».

Nous retrouvons la même ambiguïté dans la façon dont
Hume prétend résoudre le problème moral. Le bien et le
mal ne sont pas des relations objectives. Nous ne pouvons
pas fonder nos devoirs sur quelque principe qui existerait
en soi. Par exemple la volonté divine telle qu’elle peut
être connue par la philosophie ne peut pas justifier la
condamnation du suicide. On dit communément que la
Providence ne permet pas aux hommes de disposer de
leur vie. Mais c’est là bien mal argumenter. En effet, tout
ce qui se passe dans l’univers, sans exception, « peut être
%36
%{\it }
appelé l’action du Tout-Puissant », car tout procède de
lui par le truchement des pouvoirs dont il a doté ses créatures. 
« Une maison qui s'écroule d’elle-même n’est ni
plus ni moins détruite par la Providence que si elle est
démolie par les mains de l’homme. » Rien n’échappe aux
lois de la nature. Les conséquences de nos jugements,
de nos passions, les tempêtes, les inondations, les épidémies 
sont également des effets de la Providence. L’homme
qui se tue ne trouble donc pas l’ordre de l'univers, mais
l’exprime. Si la disposition de la vie humaine était réservée
à l’action particulière du Créateur,  « il serait aussi criminel
de sauver une vie que d’en détruire une ». En se plaçant
au point de vue ontologique, on ne peut donc pas plus
condamner le suicide que le meurtre. Cela Hume ne le
dit pas mais c’est une conséquence logique et l’{\it Essai
sur le suicide} fait penser aux raisonnements du marquis
de Sade qui justifie son immoralisme par un « chantage ontologique » du même style.

La « raison » laïcisée, chère à certains philosophes du
{\footnotesize XVIII}$^\text{e}$ siècle, ne peut pas davantage fonder la morale.
Car la pure raison, si elle éclaire la relation de moyen à
fin, ne fournit elle-même aucune fin : «Si une passion ne
se fonde pas sur une fausse supposition et si elle ne choisit
pas des moyens impropres à atteindre ses fins, l’entendement 
ne peut ni la justifier ni la condamner. Il n’est
pas contraire à la raison que je préfère la destruction du
monde entier à une égratignure de mon petit doigt. »
Les moralistes rationalistes (à la manière de Clarke)
passent subrepticement des jugements de réalité (qu’il
s’agisse d’affirmations métaphysiques ou d’autres propositions) 
aux jugements de valeur. Lisons leurs ouvrages :
Tout à coup, dit Hume, au lieu de la copule {\it est}, qui était
habituelle dans leurs propositions, « nous ne rencontrons
que des propositions où la liaison est établie par {\it doit} ou
ne doit pas ». Et ce tour de passe-passe n est aucunement justifié.

Où donc se trouve la source du jugement moral ?
%37
%{\it }
Uniquement dans nos sentiments, dans nos réactions subjectives. 
« La morale est plus proprement sentie que jugée. »
Une action, un sentiment, un caractère seront déclarés
«vicieux » ou « vertueux ». Pourquoi ? Parce que la vue
ou le récit de tel acte provoque en nous « un plaisir ou un
malaise d’un genre particulier ». Nous éprouvons un sentiment 
d’approbation ou de désapprobation. Le jugement
éthique, comme le jugement esthétique, repose sur des
tendances spécifiques de la nature humaine. La reconnaissance 
de la radicale subjectivité du jugement éthique, de
son relativisme, n’est nullement, comme M. Leroy le
fait remarquer, une « liquidation de la morale ». La réduction 
psychologique de l’idée de cause ne nous empêchait
pas de croire à la causalité. La réduction psychologique de
Vidée de bien ou de mal ne peut guère nous troubler ;
d’abord parce que « nous ne pouvons pas plus changer nos
sentiments personnels que le mouvement des cieux »,
ensuite parce que l’uniformité de la nature humaine
(thème très cher à Hume, nous le savons) confère à nos
sentiments éthiques une universalité remarquable. Remarquons 
ici que Hume n’ignore nullement la diversité des
mœurs et des coutumes, la multiplicité des civilisations
dans l’espace, leur évolution dans le temps. Ce n’est
pas là une découverte des sociologues ultérieurs. Montaigne — que 
Hume avait lu — avait largement insisté
sur ce thème et Hume le reprend lui-même dans le {\it Dialogue
avec Palamède}. Sur bien des articles, la morale des anciens
Grecs et des Européens du {\footnotesize XVIII}$^\text{e}$ siècle paraît s’opposer.
Mais il y a des sentiments permanents (nul peuple n’a
jamais prôné la lâcheté ou la déloyauté). Et les divergences
éthiques représentent le plus souvent des applications
diverses, dans des situations différentes, de principes
uniformes de la nature humaine. La pratique de l’exposition 
des nouveau-nés, qui révolte la conscience moderne,
n’excluait pas des sentiments de compassion : la mort
passait alors pour moins redoutable que la misère.

Cependant Hume n’ignore pas qu’un fondement 
%38
%{\it }
purement empirique de la morale est plutôt une explication
psychologique qu’une justification philosophique. Si je
parle à un malhonnèête homme, reconnaît-il, du « contentement 
intérieur qui résulte des actions louables et
humaines, du plaisir délicat de l’amour et de l’amitié
désintéressés, il pourrait me répondre que ce sont là des
plaisirs pour qui est capable de les goûter, mais que, pour
sa part, il se trouve d’un tour d’esprit et de dispositions
tout à fait autres. Ma philosophie n’apporte aucun remède
à un tel cas ». Les héros du marquis de Sade ne parleront
pas autrement de la morale et du sentiment. Et si Hume
s'efforce toujours, à partir de considérants sceptiques,
d'aboutir à des conclusions rassurantes, on peut se demander 
dans quelle mesure, en voulant persuader le lecteur,
il parvient à se convaincre lui-même. Ayant écrit : « Quand
on déclare une action vicieuse on veut simplement dire,
qu'étant donné la constitution {\it particulière} de sa nature,
on a un sentiment de blâme en assistant à la dite action »,
Hume, dans une lettre du 16 mars 1740, soumet son texte
à Hutcheson, et lui demande s’il n’est pas dangereux de
s’exprimer en ces termes. Il est curieux de noter que, dans
le livre III du {\it Traité de la nature humaine} qui paraît chez
l'éditeur Thomas Longman à la fin de 1740, le mot « particulière » 
disparaît de ce texte. Hume a littéralement ici
effacé le problème, ne pouvant le résoudre.

On retrouverait les mêmes oscillations dans la pensée
politique de Hume. Sa réflexion critique le porterait vers
des solutions révolutionnaires ou tout au moins libérales.
Pas plus pour lui que pour Pascal les « grandeurs d’établissement » 
ne sont respectables. Il est bien « rare de
rencontrer une race de rois ou une forme de république qui
ne se fonde pas primitivement sur l’usurpation et la
violence ». La science historique disqualifie le traditionalisme 
tout comme la psychologie démystifie la causalité.
Et pourtant, en pratique, « il n’y a pas de maxime plus
conforme à la prudence que de se soumettre paisiblement
au gouvernement que nous trouvons établi dans le pays
%39
%{\it }
où il nous arrive de vivre ». La critique spéculative débouche
sur le conservatisme pratique. Ici encore la pensée est une
chose, la vie en est une autre. Pensons comme les whigs
mais votons pour les tories !
%%%%%%%%%%%%%%%%%%%%%%%%%%%%%%%%%%%%%%%%%%%%%%%%%%%%%%%%%%%%%%%%%%%%%%%%%%%

%

%%%%%%%%%%%%%%%%%%%%%
\section{Le problème religieux}
%%%%%%%%%%%%%%%%%%%%%
%{\it }

Nous retrouvons, en étudiant les réflexions de Hume
sur la religion, la distinction fondamentale — mise en
lumière dans l’étude de la causalité — entre la question
du fondement et celle de l’origine. Cest un problème de
savoir si la notion de connexion nécessaire est ou non
fondée en raison. C’est un autre problème de savoir pourquoi,
en fait, nous croyons à la connexion nécessaire ;
ce deuxième problème substitue à la recherche philosophique
du fondement l’enquête psychologique sur l’origine.
Les deux questions se posent également à propos de l’affirmation
religieuse. Au tout début de sa dissertation sur
l’{\it Histoire naturelle de la religion}, Hume déclare que toute
enquête sur la religion doit s'attacher à résoudre deux
problèmes, celui du fondement de la religion {\it (its foundation in reason)}
et celui de ses sources psychologiques {\it (its origin in human nature)}.
Le premier problème porte sur
les preuves philosophiques de l’existence de Dieu (notamment,
pour Hume, la preuve de Dieu à partir des merveilles de la nature),
le second sur les origines psychologiques de la croyance religieuse.
Le premier problème sera
traité dans les {\it Dialogues sur la religion naturelle}, le second
fait l'objet de l’{\it Histoire naturelle de la religion}. C’est peu
de dire que les deux problèmes sont distincts, que les
motifs pour lesquels, selon les philosophes, nous devrions
croire sont différents des mobiles pour lesquels la plupart
des hommes croient effectivement. En fait, les motifs
philosophiques théoriques et les mobiles psychologiques
réels sont radicalement opposés. Le philosophe, pour
prouver Dieu, invoque l’ordre de la nature, l’enchaînement
%40
%{\it }
des phénomènes, admirable dans sa régularité, l’heureuse
harmonie des organes dans les êtres vivants. Tout cela
laisse indifférents les hommes vulgaires : « Plus la nature
est régulière et uniforme, moins ils ont envie de la creuser. »
L’accoutumance dissimule l'énigme du monde, étouffe
toute inquiétude philosophique. C'est au contraire le
désordre, les accidents, les catastrophes imprévues qui
font que les hommes interrogent les puissances célestes.
Les joueurs, les mariniers, qui à tout instant dépendent du
hasard, sont les hommes les plus dévots. Les maladies,
le malheur, la mort nous font « plus aisément fléchir les
genoux » que les événements heureux : Bref, «ce qui, chez
un bon esprit, fait une des plus fortes objections contre
l'existence de l’Étre suprême est pour l’homme du commun
le seul argument pour Lui ». 

La croyance en Dieu comme toutes les autres croyances,
prend sa source dans la « nature humaine », dans nos
passions de crainte et d'espérance. Ce « psychologisme »
fait de l'{\it Histoire naturelle de la religion} un ouvrage très
original à son époque. Non seulement pour Hume le
polythéisme est le point de départ chez tous les peuples
de la croyance religieuse (et non pas la dégradation d’un
prétendu monothéisme primitif); mais encore ce polythéisme
primitif ne relève aucunement de la psychologie
intellectualiste qu’on lui appliquait communément ; les
dieux, bienveillants ou maléfiques, des peuples primitifs,
ne sont pas tant des principes explicatifs des choses que
la projection, dans le ciel, de nos angoisses et de nos espoirs.

A ce polythéisme d’origine psychologique s'oppose
cependant le déisme rationnel du philosophe. Tandis que
le peuple « ne monte jamais aux cieux par la contemplation »,
l’idée d’une « Cause intelligente » de l'Univers
frappe le philosophe « comme une évidence qui porte
conviction ». Hume l’affirme à plusieurs reprises dans
l'{\it Histoire naturelle de la religion}. Faut-il voir dans ces
affirmations une précaution verbale, une habileté de Hume
qui espère ainsi faire passer l'audace de ses analyses psychologiques
%41
%{\it }
et historiques de la religion populaire ? Ou
bien Hume est-il un déiste sincère ?

Pour tenter de répondre à cette question, il faut examiner
de près les {\it Dialogues sur la religion naturelle}. L'ouvrage
avait paru très destructeur aux rares amis auxquels le
philosophe l’avait confié. Blair pensait qu’il eût mieux
valu le détruire, Adam Smith refusa, comme on sait, la
responsabilité de sa publication posthume. L'ouvrage se
présente comme le récit par le jeune Pamphile — élève
du philosophe Cléanthe — d’une suite de discussions entre
trois personnages : Cléanthe, déiste résolu, partisan de la
preuve de Dieu par la finalité du monde ; Philon le sceptique,
Déméa, mystique antirationaliste. Il est manifeste,
bien que les historiens de Hume ne l’aient pas toujours
remarqué, que ce livre est inspiré du {\it De Natura Deorum}
de Cicéron, longue discussion de Cotta, de Velleius et de
Balbus sur les conceptions épicurienne et stoïcienne de la
divinité. Hume emprunte à Cicéron non seulement la
forme du dialogue mais beaucoup d’idées ; Cicéron avant
Hume déclare : « Sur la nature des Dieux on varie, personne ne
nie leur existence. » Le héros déiste s’appelle
Cléanthe comme le maître stoïcien. Comme lui il rejette
le hasard épicurien et insiste sur l’évidence de la providence.

Quel est dans ces dialogues le porte-parole de Hume ?
Pamphile, à la fin de l’ouvrage, dit que Cléanthe a raison.
Mais Pamphile n’est pas Hume et c’est le jeune disciple
de Cléanthe. Il conclut tout naturellement en faveur de
son maître ! Hume lui-même s’est longuement expliqué
sur le sens de son ouvrage dans une longue lettre de 1751
à son ami Gilbert Elliot of Minto. Il lui confie qu'avant
l’âge de vingt ans il avait écrit un gros manuscrit sur le
problème religieux, brûlé par la suite et dont la substance
a dû passer dans les {\it Dialogues}. Il « commençait par une
anxieuse recherche des arguments susceptibles de confirmer
l’opinion commune. Puis des doutes surgissaient,
étaient dissipés. revenaient encore ». Or la forme dialoguée
%42
%{\it }
est celle qui convient le mieux à une pensée hésitante.
Hume écrit à son ami qu’il a « voulu éviter cette vulgaire
erreur qui consiste à ne mettre que des absurdités dans
la bouche des adversaires ». Il ajoute, il est vrai, que
Cléanthe est le héros du livre, et il prie Gilbert Elliot de
lui communiquer, après lecture du manuscrit des {\it Dialogues},
« toutes les pensées qui vous viendront pour fortifier la position
de Cléanthe ». Mais,c’est là une remarque
{\it ad hominem} car précisément les opinions philosophiques
d’Elliot sont telles que si on jouait les {\it Dialogues}, ce dernier
« pourrait tenir fort bien le rôle de Cléanthe ». Et Hume
ajoute, ce qui est capital : « Pour mon compte, je me
réserverai le rôle de Philon que vous le reconnaissez,
je pourrai tenir avec assez de naturel. » Et Hume aura
beau écrire plus tard, à Balfour, qu’il « s’est efforcé de
réfuter le sceptique Philon » cela n’empêche pas qu'il
donne dans les {\it Dialogues} au philosophe sceptique la
part du lion. Comme le remarque J. Y. T. Greig : « Les
discours de Cléanthe couvrent trente-cinq pages, ceux de
Philon, qui doit être réfuté et réduit au silence, couvrent
cent quatorze pages. »

La preuve de Dieu par les causes finales, cependant, ne
manque pas d’impressionner Philon-Hume. C’est la preuve
favorite des déistes du {\footnotesize XVIII}$^\text{e}$ siècle. Leur optimisme
naturaliste s’y reflète complaisamment. Kant lui-même,
étranger à tout naturalisme, trouvera encore cet argument
« respectable ». Hume l’avait découvert très tôt dans le
{\it De Natura Deorum}. Le philosophe stoïcien auquel Cicéron
donne la parole conteste aux épicuriens que « ce monde
si riche et si beau » puisse résulter d’un concours fortuit
d’atomes : « Si l’on croit cela, pourquoi ne penserait-on
pas également que, en jetant d’une manière quelconque
en quantité innombrable les vingt et une lettres, il pourrait
résulter de ces lettres, jetées sur le sol, les {\it Annales}
d’Ennius telles qu’on pourrait les lire ensuite ? Je ne sais
si le hasard pourrait seulement venir à bout même d’une
seule ligne. » Le Cléanthe de Hume reprend inlassablement
%43
%{\it }
l’argument. La structure d’une fleur, l’anatomie du
moindre animal offrent maintes preuves plus fortes d’un
dessein originel qu’une demeure construite par un architecte ou qu’un livre de Tite-Live ou de Tacite. Tout
comme un livre est produit par l'intelligence humaine,
le « soigneux ajustement des causes finales » est un langage
dans lequel nous pouvons aisément déchiffrer l'intention
d’un Créateur souverainement intelligent.

Mais Philon, un instant déconcerté, soumet la preuve à
une critique serrée. Le raisonnement analogique qui nous
fait expliquer la production du monde matériel par une
Intelligence suprême, analogue à l'intelligence humaine
quoique infiniment supérieure, manque de rigueur. La
nature dans son ensemble est malgré tout trop différente
des produits de l’industrie humaine pour que ce raisonnement
soit décisif, « Le monde, dans son ensemble,
ressemble manifestement plus à un animal où à un végétal
qu’il ne fait à une montre et à un métier à tricoter. » Et
que gagne-t-on à expliquer l’ordre de la matière par
l'œuvre d’un esprit ? Sans doute l'organisation des parties
d’une maison témoigne-t-elle de l'intelligence de l’architecte.
Mais la façon dont les idées s’organisent dans l’esprit
de l’architecte lui-même demeure une énigme. L’intelligence de
l’homme suppose l’organisation ; elle n’en est
pas la clef. L’ordre n’est pas expliqué par un dessein,
mais un dessein suppose lui-même un ordre. On dira que
Dieu est précisément l'explication ultime de cette énigme.
Mais inventer un Dieu revient à redoubler l’énigme plutôt
qu’à la résoudre. Si l’ordre du monde s’explique par
l’intelligence divine, il faudrait expliquer comment les idées
se disposent en ordre dans cette intelligence suprême.
Et si nous alléguons que « les idées qui composent la raison
de l’Être Suprême se disposent en ordre d’elles-mêmes et
par leur propre nature », nous pouvons aussi bien faire
l’économie de ce Dieu et dire que « les parties du monde
matériel se disposent en ordre d’elles-mêmes et par leur
propre nature ». Il n’est pas plus clair d'expliquer la
%44
%{\it }
matière par l'esprit que l’esprit par la matière :  « la
matière peut contenir la source de l’ordre originellement
en elle-même aussi bien que l'esprit. » 

Et puis à quoi bon s’émerveiller sur l'extraordinaire
adaptation des organes chez les animaux et chez les
plantes ? « Je voudrais bien savoir comment un animal
pourrait subsister à moins que ses parties ne fussent
adaptées. » Ici Hume annonce Auguste Comte et Darwin.
Le principe des « conditions d’existence » se substitue au
principe de finalité. La sélection des plus aptes est en
quelque sorte automatique et ne résulte pas d’un dessein :
« Bien des mondes auront pu être gâchés durant une éternité
avant que ce système ne fût mis au jour. » Cependant
Philon reconnaît que sa réfutation est subtile et laborieuse.
Devant la simplicité de l’argument téléologique, il avoue
que « toutes les objections paraissent arguties et sophismes
purs ». 

Accordons à Cléanthe que l’univers témoigne d'une
finalité. Seulement (et c’est ici que Philon reprend l’avantage)
cette finalité n’a rien à voir avec ce qu'on appelle
une providence. Les soigneux artifices de la nature semblent
n’avoir pour but que de rendre « plus amère la vie de tout
être vivant ». Les espèces animales et végétales ont été
dotées d’armes offensives et défensives parfaitement
adaptées à la lutte de tous contre tous. La création est le
théâtre d’un perpétuel carnage : « Tout animal est entouré
d’ennemis qui recherchent incessamment sa misère et sa
destruction. » L’homme lui-même est sans cesse menacé
par les maladies, les inondations, les tremblements de
terre. Les pages très éloquentes que Hume, dans la
Dixième partie des {\it Dialogues}, consacre au problème du
mal nous font penser à cette remarque de Léon Brunscvicg. Si
le Dieu créateur « est un Dieu artiste, n’est-ce
pas au sens néronien du mot » ? Toutes les justifications
philosophiques du mal sont dérisoires ; on dit que la
douleur est un avertissement; mais une véritable providence eût
donné aux êtres vivants des signaux moins
%45
%{\it }
catastrophiques ; on prétend que le mal résulte inévitablement,
à titre d’accident, des lois générales de la création ; mais
une providence bienveillante n’eût pas dédaigné
d’agir par des volontés particulières. Et pourquoi les
forces des êtres vivants sont-elles si étroitement limitées ?
Enfin à supposer que les vents, les pluies, la distribution
de la chaleur et du froid résultent d’une volonté bienveillante,
pourquoi ces « ressorts » sont-ils si souvent « déréglés »,
aboutissant à la destruction de ces vies qu’ils seraient
censés protéger ?

Si Dieu est bon il n’est pas tout-puissant, s’il est tout-puissant
on ne peut parler de sa bonté. Le vieil argument
d’Épicure demeure irréfutable. Sans doute, si nous avions
par ailleurs des preuves certaines de l’existence d’un
créateur souverainement bon, on pourrait s’exercer à
montrer que l’existence du mal dans l’univers n’est pas
{\it incompatible} avec la Providence. Mais on ne saurait, de
toute évidence {\it inférer} la bonté de Dieu du spectacle de
l’Univers ! Ce spectacle, à le considérer sans préjugé,
« n’éveille pas d’autre idée que celle d’une nature aveugle,
imprégnée par un grand principe vivifiant et laissant
tomber de son giron, sans discernement ni soin maternel,
ses enfants estropiés et avortés ».

De toute façon la critique du principe de causalité,
thème fondamental de la philosophie de Hume, n’a-t-elle
pas d’avance dévalué toutes les preuves de l'existence de
Dieu ? Nous avons vu que l’inférence causale n’est qu’une
habitude née de notre expérience passée. Or Dieu n’est
pas une donnée de l’expérience. C’est une cause inventée,
non une Cause observée. Hume rappelle, dans ses {\it Dialogues},
la découverte fondamentale des ouvrages antérieurs :
«Il n’y a pas d’être dont la non-existence implique contradiction ? Par conséquent il n’y a pas d’être dont l’existence
soit démontrable. »

Si nous nous en tenions là, il faudrait convenir que le
scepticisme de Hume est tout proche de l’athéisme. Mais
les discussions entre Philon et Cléanthe ne doivent pas
%46
%{\it }
nous faire oublier les étranges relations entre Philon et
Déméa. Le sceptique et le « mystique » sont toujours
d’accord pour critiquer le rationalisme et l'optimisme de
Cléanthe. Sans doute Hume semble-t-il vouloir montrer,
non sans humour, que le mystique est  « un athée qui
s’ignore ». Parler comme Déméa de « la nature adorablement mystérieuse de Dieu », ou dire comme Philon que
Dieu est inconnaissable, ne revient-il pas au même
Mais dans les dernières pages du livre nous voyons que
c’est Philon qui, tout au contraire, rejoint le point de vue
de Déméa. Le scepticisme, en humiliant le dogmatisme
orgueilleux de la raison naturelle, ne nous prépare-t-il 
pas à accueillir la révélation ? Le secpticisme philosophique
n'est-il pas une propédeutique à la foi ? « Une personne
pénétrée d’un juste sentiment des imperfections de la
raison naturelle, dit Philon, volera à la vérité révélée
avec la plus grande avidité... Être un sceptique philosophique
c’est, chez un homme lettré, le premier pas et le
plus essentiel menant à être un vrai chrétien, un croyant »
N'est-ce pas en ce sens que Hume entendait le mot de
Bacon qu’il cite volontiers : « Un peu de philosophie
éloigne de la religion, beaucoup de philosophie y ramène ? »
On ne peut s’empêcher ici de penser à Pascal : « Le pyrrhonisme
sert à la religion. » Songeons, par exemple, que
tandis que Kant écrivait en 1780 sa {\it Critique de la raison
pure}, Hamann lui envoya un travail sur les « Dialogues » de
Hume dans l’espoir que la critique du Dieu des philosophes
disposerait Kant à accueillir le Dieu de la révélation chrétienne.

On objectera peut-être que l’appel à la révélation n’est
chez Hume, que persiflage et pirouette. Hume est sévère
à l'égard du christianisme. Pensant à sa propre enfance,
il rappelle les critiques que Machiavel — précurseur de
Nietzsche sur ce point — adressait à l’éducation chrétienne qui,
en insistant sur les mortifications, la pénitence,
l’obéissance, décourage l’énergie et débilite la {\it virtu}. Il
déteste les prêtres — au moins autant que les Anglais —
%47
%{\it }
(ce qui, pour un Écossais bon teint, n’est pas peu dire).
Dans son {\it Histoire d'Angleterre}, il condamne l’autoritarisme
clérical et le fanatisme. Et ce qui est plus important, dans
Son {\it Essai sur les miracles}, il raille la foi en la résurrection
du Christ, sous un déguisement facile à percer (« Supposez
que tous les historiens qui traitent de l’Angleterre s’accordent sur ce
que le 1$^\text{er}$ janvier 1600 la reine Elisabeth
mourut... qu'après un mois d’inhumation elle reparut,
réoccupa le trône et gouverna l’Angleterre pendant trois
ans »). Hume rejette les miracles au nom de la causalité
physique, tout en expliquant la croyance au miracle par
les passions de la nature humaine, par la causalité psychologique.
Mais précisément, quand on sait ce que Hume
pense du principe de causalité (simple habitude mentale
sans valeur rationnelle), on peut considérer que sa critique
du miracle, qui reste purement psychologique (humainement il n’est
guère possible de croire au miracle), n’est
Pas son dernier mot. Une résurrection, disait Pascal,
est ni plus ni moins énigmatique qu’une naissance ;
seule l’habitude nous. rend l’une croyable, l’autre incroyable,
« populaire façon d’en juger ». Et Hume lui-même, à la fin de
l’{\it Essai sur les miracles}, réserve la possibilité de ce miracle
intérieur qu’est la foi du croyant,
« miracle continu dans sa propre personne qui lui donne
une détermination à croire ce qui est le plus contraire
à la coutume ». L’essai sur l'{\it Immortalité de l’âme} se termine
de la même façon. Et Philon, à la fin des {\it Dialogues}, s’écrie :
« Plaise au Ciel d’alléger notre profonde ignorance en
offrant à l’humanité quelque révélation particulière. »
Persiflage pur et simple, dira-t-on, mais qui atteint la
philosophie de Hume lui-même Car un empiriste conséquent
ne peut soulever d’objection contre ce fait, contre cette
expérience intime vécue d’une révélation particulière !
Assurément, Hume ne prétend nulle part avoir reçu
pareille révélation. Du moins y a-t-il là une issue possible,
une porte que son impitoyable critique laisse entrouverte.
%{\it }
%%%%%%%%%%%%%%%%%%%%%%%%%%%%%%%%%%%%%%%%%%%%%%%%%%%%%%%%%%%%%%%%%%%%%%%%%%%

%

%
\vfill
\begin{center}
\includegraphics[scale=0.7]{./presentation/chatDeDos}
\end{center}
\vfill
%
\chapter{Extraits}
%
\input{./extraits/.tex}
%

%
%
\begin{appendix}
%

%%%%%%%%%%%%%%%%%%%%%
\chapter{La vie}
%%%%%%%%%%%%%%%%%%%%%
%{\it }{\oe}

David Hume est né le 26 août 1711 à Édimbourg
où son pére exerçait la profession d’avocat. Ce
dernier étant mort en 1714, Mrs Hume se retira
avec ses trois enfants John, Katherine et David
dans la propriété familiale de Ninewells, le domaine
des « Neuf Sources » situé dans la pittoresque campagne
(avec ses falaises, ses ruisseaux et ses bois)
du Berwickshire. L’oncle de David, pasteur du
village voisin de Chirnside, dirigea sa toute première
éducation. L’enseignement religieux que reçut le
jeune David semble avoir été particulièrement
austère et maladroit (le Révérend George Hume
se plaisait, dans ses sermons, à humilier publiquement
les jeunes filles dont la grossesse révélait les
péchés charnels). L’antipathie précoce de Hume
pour le christianisme vient en partie de là.

Cependant le petit David échappa assez rapidement
à cette atmosphère déprimante. Élève dès
l'âge de onze ans du collège d’Édimbourg (renommé
à juste titre et qui deviendra plus tard université),
il se trouve dans une ambiance intellectuelle beaucoup 
plus stimulante. Il y écoute les cours de « philosophie
naturelle », c’est-a-dire de physique, de
Robert Stewart (disciple de Newton aprés avoir été
cartésien) et se souviendra certainement de ses
%6
%{\it }{\oe}
leçons lorsqu’il rêvera d’appliquer la méthode expérimentale
à la morale et à la métaphysique. Mais
la formation de Hume au collège fut essentiellement
littéraire. C’est de cette époque que date son goût
de Virgile et de Cicéron. Ce sont les textes (notamment
le {\it De Natura Deorum}) où Cicéron résume les
débats philosophiques des Stoïciens et des Epicuriens
qui découvrent à Hume le monde des discussions
métaphysiques. Revenu à Ninewells dès sa quinzième
année, le jeune David se livre avec passion à la
lecture des anciens et des modernes. Il dévore
Montaigne, Bacon, Malebranche, Bayle, mais aussi
Milton, Pope, Swift, Shaftesbury. À l’âge de vingt
ans, il a déja rempli un gros cahier de réflexions sur
le problème religieux, sur la psychologie, sur l’histoire.
Cette activité intellectuelle bouillonnante, les
leçons de scepticisme qu’il tire de lectures aussi
diverses, un conflit avec sa famille qui voudrait,
malgré lui, l'orienter vers des études juridiques
provoquent une crise de dépression passagère dont
nous trouvons le témoignage dans un curieux brouillon
de lettre à un médecin célèbre (qui n’est pas
comme on l’a cru George Cheyne, mais le D$^\text{r}$ Arbuthnot).
David qui n’a hérité de son pére que d’une
toute petite rente doit de toute urgence prendre
un état. Après un bref essai dans le commerce (au
service d’un marchand de Bristol), David Hume
décide de ne plus résister à sa vocation : Il sera
philosophe et homme de lettres, et il entend conquérir
la gloire. Pour pouvoir subsister, il se rend en France
(où la vie est à l’époque beaucoup moins chère),
s’installe à Reims en 1734 à l'hôtel du {\it Perroquet
vert}, puis à La Flèche en Anjou, tout près du collège
%7
%{\it }{\oe}
de Jésuites ou Descartes fut élève. Il y rédige, à
peine âgé de vingt-trois ans, son chef-d’{\oe}uvre :
le {\it Traité de la nature humaine}.

Revenu à Londres en 1737, il a la chance de trouver
un éditeur, et la prudence (ou si l’on veut la faiblesse)
de supprimer les chapitres sur la religion
(il espère la protection de l’évèque Butler). Les deux
premiers livres du traité ainsi publiés « tombérent
mort-nés de la presse », raconta plus tard Hume dans
sa courte {\it Autobiographie}. Ce n’est pas tout à fait
vrai. En fait l’{\oe}uvre intéressa quelques critiques,
mais n’atteignit pas le grand public (qui seul donne
la notoriété). A cette époque, Hume entre en relations
avec Hutcheson, professeur à Glasgow, qui
lui présente son jeune étudiant Adam Smith (qui
restera toujours l’ami de Hume) et lui trouve un
éditeur pour les deux livres suivants du {\it Traité de
la nature humaine} dont le succès n’est pas plus
grand. Hume cependant ne doute pas de sa valeur.
Son échec vient de la présentation trop lourde et
trop savante de sa philosophie et non pas du fond
({\it more from the manner than the matter} dira l'{\it Autobiographie}).
Hume décide alors d’écrire des essais
courts et brillants et publie en 1741 a Édimbourg
ses {\it Essais moraux et politiques} (humilié par ses
insuccés il présente d’ailleurs cet ouvrage comme son
premier livre !). Cette fois les lecteurs sont nombreux,
et Hume croit pouvoir présenter sa candidature à
la chaire de philosophie morale de l’Université de
Glasgow. L’opposition des chrétiens empêche sa
nomination. En 1746, Hume devient le secrétaire
particulier du général Saint-Clair, un Écossais qui
est son parent éloigné, et l’accompagne dans une
%8
%{\it }{\oe}
mission diplomatique à Vienne et à Turin. Pendant
son voyage paraissent ses {\it Essais philosophiques sur
l'entendement humain} (plus tard {\it Enquête sur l’entendement
humain}), qui reprennent dans un style nouveau
les deux premiers livres du {\it Traité de la nature
humaine}. Cette fois les chapitres sur le miracle et
sur la providence particulière paraissent avec le
reste (1748). A son retour c’est le troisième livre du
Traité que Hume reprend avec l’{\it Enquête sur les
principes de la morale} (1751). Dès lors la notoriété
de Hume s’affirme. Il entre en relations épistolaires
avec Montesquieu à propos de l’{\it Esprit des lois}.
Et s'il échoue une nouvelle fois (1751) dans sa candidature
à l’Université de Glasgow, il devient conservateur
de la bibliothéque de la « Faculté des Avocats »
à Edimbourg. Il trouve là tous les documents
nécessaires pour écrire de 1754 à 1761 sa monumentale
{\it Histoire d’ Angleterre}, de Jules César à Jacques II.
Le premier volume qui traite des régnes de Jacques
I$^\text{er}$ et de Charles I$^\text{er}$ déclenche un petit scandale
dans les milieux religieux, et Hume n’améliore pas
son cas auprès des chrétiens en publiant ses {\it Quatre
dissertations} (dont son {\it Essai sur l'histoire naturelle
de la religion} et son {\it Essai sur le suicide}). Il est vrai
qu'il s’empresse de retirer l’{\it Essai sur le suicide}
ainsi qu’un {\it Essai sur l’immortalité de l’âme} pour les
remplacer par un {\it Essai sur la règle du goût}. Hume
subit aussi avec patience les tracasseries des avocats
d’Édimbourg qui lui reprochent d’avoir acheté
pour la bibliothéque les {\it Contes} de La Fontaine, et
l'{\it Histoire amoureuse des Gaules} de Bussy-Rabutin !

C’est en France que Hume devait connaître la
gloire. A l’appel de lord Hertford, ambassadeur
%9
%{\it }{\oe}
d’Angleterre à Paris, il va exercer de 1763 à 1766
les fonctions de secrétaire d’ambassade. Il sera
même quelque temps — lorsque lord Hertford
est rappelé et en attendant son successeur — « chargé
d’affaires », c’est-à-dire en fait ambassadeur, grâce
à la puissante protection de la comtesse de Boufflers.

A Edimbourg Hume inquiète, à Londres il n’est
qu’un Écossais, qu’un intellectuel provincial. A
Paris, le petit monde des philosophes qui a lu ses
Essais, qui connaît l’opposition de Hume à la
« superstition » et au « fanatisme » le tient pour un
philosophe de premier plan. Il deviendra très vite
l'ami de d’Alembert, de Diderot, d’Helvétius, du
baron d’Holbach. Certes, officiellement, Hume est
déiste, à un diner du baron d’Holbach il confesse
même n’avoir jamais rencontré d’athée (« Regardez
autour de vous, répond le baron : Il y en a quinze
autour de cette table ! »). En fait ses positions anti-religieuses,
son Essai sur les miracles le rendent
sympathique aux encyclopédistes qui le tiendront
désormais pour un « frére ». A Paris, c’est un véritable
triomphe. Les plus grandes dames s’arrachent ce
quinquagénaire bedonnant. C’est la duchesse de
La Vallière qui tient à le voir dès son arrivée à Paris,
avant même qu’il ait pu changer de costume !
C’est Mme Du Deffand, Mme Geoffrin, Mlle de Lespinasse
qui le fêtent dans leurs « salons ». même
l'apparence physique de Hume, assez ingrate (il
est obèse, son visage empâté est peu expressif)
qui lui avait valu naguére, tandis qu’il voyageait
en Italie avec le général Saint-Clair, les railleries
du jeune James Caulfield, futur lord Charlemont,
et quelques déboires amoureux, est maintenant
%10
%{\it }{\oe}
trouvée sympathique. Son visage un peu lourd, son
fort accent écossais donnent au grand philosophe
un air débonnaire, une simplicité de bon aloi
(Mme Du Deffand l'appelle « mon cher paysan »,
Mme Geoffrin « mon gros drôle, mon gros coquin »).
Hume vole de succès en succès, écrit à son ami le
D$^\text{r}$ Robertson : « Je ne mange que de l’ambroisie,
je ne bois que du nectar, je ne respire que de l’encens,
je ne marche que sur des fleurs. »

Marie-Charlotte Hippolyte de Campet de Saujeon,
comtesse de Boufflers et maîtresse en titre du prince
de Conti, avait dès 1761 écrit à Hume pour lui dire
toute son admiration pour sa philosophie « sublime ».
Elle avait tenté de le rencontrer au cours d’un voyage
à Londres, sans succés car Hume n’avait pas voulu
quitter Édimbourg. A Paris elle l'invite tout de
suite à ses lundis, puis à ses vendredis plus intimes.
Il semble que cette jolie femme de trente-cinq ans
ait été quelque peu amoureuse du gros philosophe
vieillissant. Hume la considéra toujours comme une
amie très chère (le philosophe, cinq jours avant sa
mort, lui écrit encore le 20 août 1776 pour lui
annoncer qu’il se sent perdu et lui dire une dernière
fois son « affection » et son « respect »), mais il ne
paraît pas que leurs relations aient jamais été plus
intimes. Nous connaissons mal la vie intime de Hume,
mais il semble que le philosophe, resté célibataire
(à Edimbourg, sa s{\oe}ur Katherine tenait son ménage)
se soit, à cause d’anciens déboires, ou par souci de
préserver son indépendance et son travail, toujours
méfié des passions amoureuses.

En 1766, le nouvel ambassadeur d’Angleterre, le
duc de Richmond, arrive à Paris et Hume repart
% 11
%{\it }{\oe}
en Angleterre. C’est ici que se situe un des épisodes
les plus mal éclaircis de la vie de Hume. Dès ses
lettres de 1761 la comtesse de Boufflers avait intéressé
Hume à Jean-Jacques Rousseau. Pendant le
séjour à Paris, la marquise de Verdelin demande au
philosophe écossais de chercher pour Jean-Jacques
un refuge en Angleterre. Celui-ci, proscrit de Genéve,
sa ville natale, interdit de séjour en France, persécuté
par les habitants de Môtiers-Travers, en butte à la
haine des encyclopédistes qui le tiennent pour un
dévot, est dans une des périodes les plus critiques de
son existence. Hume, ému par les malheurs de
Jean-Jacques, et tout d’abord enthousiasmé par sa
simplicité et sa franchise (Rousseau est un Socrate
moderne, dira-t-il), part avec lui le 4 janvier 1766.
Ils arrivent à Londres le 13 ou Jean-Jacques est
fété et reçoit de Hume mille témoignages d’amitié.
Mais Rousseau qui déteste le monde et cherche la
solitude n’entend rester à Londres que jusqu’à
l'arrivée de sa servante-maîtresse Thérèse Le Vasseur.
En fait c’est seulement le 19 mars 1766 qu’il part
pour Wooton, maison de campagne dans les bois du
Derbyshire qu’un ami de Hume, Davenport, avait
mise à sa disposition.

Moins de trois mois ont suffi pour altérer la belle
amitié de Hume et de Rousseau. Désormais Rousseau
considére Hume comme un traître et comme un
malhonnête homme. Que s’est-il passé ? Il est certain
que Rousseau a toujours eu un caractère soupçonneux,
une tendance paranoiaque au délire d’interprétation
que les réelles persécutions dont il fut
victime n’ont fait qu’aggraver. Avant même d’atteindre Calais,
le premier soir du voyage, un curieux
% 12
%{\it }{\oe}
incident se produit. Hume, Rousseau et Luze, un
ami suisse qui va à Londres pour ses affaires, descendent
dans un hotel de Senlis et couchent dans une
chambre à trois lits. Jean-Jacques qui cherche en
vain le sommeil entend soudain Hume dire à plusieurs
reprises « à pleine voix » et avec « une véhémence extrême » :
« Je tiens Jean-Jacques Rousseau. »
Hume n’est-il pas l’ami intime des ennemis les plus
acharnés de Rousseau, les encyclopédistes athées,
Diderot, d’Alembert et d’Holbach ? Ne s’est-il
pas mis d’accord avec cette clique pour faire de
Rousseau en quelque sorte son prisonnier ? Dès le
départ ce rêve ou cette hallucination (car Luze
n’a pas été éveillé par ces voix véhémentes) met
Rousseau sur ses gardes.

Cela suffit-il pour que nous n’accordions aucun
crédit aux accusations formulées ultérieurement par
Jean-Jacques ? Certaines assurément sont délirantes.
A Londres, Rousseau est exaspéré par les
compliments de Hume qui a toujours à son chevet
un tome de la {\it Nouvelle Héloïse}. Pure hypocrisie,
juge Rousseau, car Hume ne peut aimer ce roman !
En réalité Hume, qui fait effectivement des réserves
sur les idées de Jean-Jacques, apprécie son roman,
le tient pour le chef-d’{\oe}uvre du philosophe français
(lettre à Blair du 25 mars 1766). Mais il y a plus
grave. Hume n’ignore pas qu’à Paris son ami Walpole
a fait à Rousseau une plaisanterie trés méchante
(écrivant et faisant publier une lettre d’invitation a
Rousseau, au nom du roi de Prusse qui lui promet à
son choix faveurs ou persécutions, puisqu’il paraît
aimer les persécutions !). Or tandis que Rousseau se
préoccupe de faire venir à Londres des papiers pour
% 13
%{\it }{\oe}
la rédaction de ses {\it Confessions}, papiers restés en
France, Hume lui propose Walpole comme commissionnaire !
Hume fait faire une enquête à la banque
Rougemont sur les vraies ressources de Jean-Jacques
qui crie toujours misère. D’autre part Hume veut
toujours se charger d’expédier et de retirer de la
poste le courrier de Jean-Jacques, et celui-ci se
plaindra que son courrier lui parvienne toujours
décacheté et grossièrement refermé. Hume répondra
plus tard que s’il tenait à se charger du courrier de
Rousseau, c’est parce que ce dernier, alléguant sa
pauvreté, ne voulait payer le port d’aucune lettre
(a l’époque, c’est le destinataire qui payait) et que
lui, David Hume, ne tenait pas a laisser trop longtemps
le courrier à la poste, mais désirait le soustraire
promptement {\it from the curiosily and indiscretion
of the clerks of Post-office} !

Quoi qu’il en soit de toutes ces accusations, il
reste que Hume porte la responsabilité d’avoir mis
la querelle sous les yeux du public en laissant éditer
à ses amis français l’{\it Exposé succinct} de son différend
avec Rousseau, livrant ainsi Jean-Jacques à de
nouvelles railleries. Ce fut peut-être en cette affaire
sa seule faute, une entreprise, écrira-t-il à Adam
Smith le 17 octobre 1767, que « l’un et l’autre nous
avons été enclins a blâmer parfois, à regretter
toujours ». Si Hume ne fut pas le traître qu’a imaginé
Rousseau, peut-être manqua-t-il ici de patience et
de générosité. Une des faiblesses de Hume est d’avoir
été toujours trés préoccupé de sa réputation. Et
il craignait que Rousseau, dans ses {\it Confessions}, ne
raconte l'histoire à sa manière (en réalité le récit de
Rousseau s’arrétera avant les événements de 1766).
% 14
%{\it }{\oe}

Hume, grace à son fidèle protecteur, lord Hertford,
devient sous-secrétaire d’État en 1767. Par un
étrange retour des choses c’est lui qui, entre autres
affaires, est chargé de régler les conflits et de décider
de l’avancement des pasteurs de cette Église
d’Écosse qui naguére avait tenté d’entraver sa
carriére !

En 1769 Hume rentre à Édimbourg désormais
riche et considéré. Il choisit de finir ses jours dans
sa ville natale, intense foyer de culture, qui est bien
plutôt que Londres la capitale intellectuelle de
l’Angleterre de ce temps. A Édimbourg, cette
« Athènes du Nord », Hume retrouve en effet des
amis éminents, Adam Smith, le juriste lord Kames,
Ferguson, et aussi des adversaires loyaux et courtois
comme le théologien George Campbell qui avait
rédigé une critique de l’{\it Essai sur les miracles}, critique
fort appréciée par Hume lui-même. Hume
s’emploie à corriger ses {\oe}uvres pour de prochaines
éditions, à répondre à son courrier, et il rejoint
fréquemment ses amis écossais au {\it Poker Club} et a
la {\it Select Society}, club philosophique et littéraire
qu’il avait fondé lui-même en 1754.

Très rapidement sa santé décline. Il souffre d’une
tumeur de l'intestin, et dès le début de 1776 il se
sait perdu. En avril il rédige son testament. Il a
dans ses papiers un ouvrage inédit commencé
dès 1751 et dont il a déja soumis à cette époque les
premiers chapitres à son ami Gilbert Elliot of Minto :
{\it Les dialogues sur la religion naturelle}. Adam Smith
n'est pas très favorable à la publication de cet
ouvrage. C’est donc le neveu de Hume qui sera
chargé de cette édition posthume. {\it Les Dialogues}
% 15
%{\it }{\oe}
paraîtront en 1779, plus de deux ans après la mort
de Hume.

Hume est mort avec la plus grande sérénité.
Dans sa brève autobiographie, rédigée le 18 avril 1776
il déclare : « Il est difficile d’être plus détaché de la
vie que je ne le suis à présent. » Le 13 août, il dit qu'il
se console d’abandonner des amis, car « {\it hélas, on
ne laisse que des mourants}, comme Ninon de Lenclos
le dit sur son lit de mort. La mort m'apparaît si
peu terrible maintenant qu’elle approche, que je
dédaigne de citer des héros et des philosophes comme
exemples de courage. Le témoignage d’une femme de
plaisir qui néanmoins était également philosophe
est suffisant ». Hume s’éteignit sans angoisse dans
l'après-midi du dimanche 25 août 1776. C’était la
veille de son soixante-cinquième anniversaire.
%%%%%%%%%%%%%%%%%%%%%%%%%%%%%%%%%%%%%%%%%%%%%%%%%%%%%%%%%%%%%%%%%%%%%%%%

%
\newpage
%

%%%%%%%%%%%%%%%%%%%%%
\chapter{L'{\oe}uvre}
%%%%%%%%%%%%%%%%%%%%% \textsc{}

%%%%%%%%%%%%%%%%%%%%%%%%%
\section{Ouvrages publiés du vivant de Hume}
%%%%%%%%%%%%%%%%%%%%%%%%%
% 
%{\it }
{\it A treatise of Human Nature} (1739-1740, 3 vol.).

{\it Essays moral and political} (3 vol., 1741-1742, 1748).

{\it Philosophical essays concerning Human Understanding}
(1748 ; à partir de 1758 le mot {\it Inquiry} remplace {\it Philosophical essays}).

{\it An Inquiry concerning the principles of Morals} (1751).

{\it Political discourses} (1752).

{\it The History of Great Britain} (1754-1757).

{\it Four Dissertations : 1. The natural history of Religion ;
II. Of the passions ; III. Of tragedy ; IV. Of the standard
of Taste (1757).} En 1755, Hume avait remis à l’éditeur
les trois premiéres de ces dissertations et une quatriéme
consacrée aux mathématiques et à la physique. Hume
la retire sur le conseil d’un ami mathématicien. Il la
remplace par deux dissertations : {\it On suicide} et {\it On the
immortality of soul.} Tandis que ces deux essais sont
déja en vente, Hume les retire et les remplace par une
seule dissertation : {\it Of the Standard of Taste. }

{\it The history of England} (1759 a 1767).

{\it Exposé succinct de la contestation qui s’est élevée entre
M. Hume et M. Rousseau} (1766).

%%%%%%%%%%%%%%%%%%%%%%%%%
\section{Ouvrages posthume}
%%%%%%%%%%%%%%%%%%%%%%%%%

{\it The life of David Hume written by himself} (1777).
{\it Two essays (On Suicide et The Immortality of the Soul)}
(1777).

%50 HUME
%{\it }
 

{\it Dialogues concerning Natural Religion} (1779).

J. H. \textsc{Burton}, {\it Life and Correspondance of David Hume,}
Edinburgh, 1846, 2 vol.

J. Y. T. \textsc{Greig}, {\it The letters of David Hume}, Oxford, 2 vol.,
1932.

R. \textsc{Klibansky} et E. C. \textsc{Mossner}, {\it }New letters of David
Hume, Oxford, 1954.

Rappelons qu’une édition anglaise classique comprend
toute l’{\oe}uvre philosophique de Hume : {\it The philosophical
works of David Hume}, éd. T. H. Green and T. H. Grose,
London, 1874-1875, 4 vol.

%%%%%%%%%%%%%%%%%%%%%%%%%
\section{Traductions françaises}
%%%%%%%%%%%%%%%%%%%%%%%%%{\it }

{\it {\oe}vres philosophiques choisies (Enquéte sur l'entendement,
Traité de la nature humaine, Dialogues de la religion
naturelle)}  traduites par Maxime \textsc{David} avec préface
de \textsc{Lévy-Bruhl}, Paris, Alcan, 1912. La traduction
Maxime \textsc{David} des {\it Dialogues sur la religion naturelle}
a été rééditée en 1964 chez J.-J. Pauvert dans la collection 
« Libertés » avec une présentation et des notes
de Clément \textsc{Rosset}.

{\it Traité de la nature humaine}, préfacé et traduit par André
\textsc{Leroy}, Editions Aubier (1$^{\text re}$ éd., 1946), 2 vol.

{\it Enquête sur l'entendement humain} (trad. \textsc{Leroy}, Aubier,
1947).

{\it Enquête sur les principes de la morale. Les quatre philosophes} 
(trad. \textsc{Leroy}, Aubier, 1947).

Il existe une traduction frangaise de 1788 des quatre
dissertations {\it (L’histoire naturelle de la. religion, Les passions, 
La tragédie, La règle du goût)}. Une nouvelle traduction 
des {\oe}uvres de Hume est en cours d’édition.

 

 
%%%%%%%%%%%%%%%%%%%%%%%%%%%%%%%%%%%%%%%%%%%%%%%%%%%%%%%%%%%%%%%%%%%%%%%%

%
\newpage
%

%%%%%%%%%%%%%%%%%%%%%
\chapter{Bibliographie}
%%%%%%%%%%%%%%%%%%%%%

%%%%%%%%%%%%%%%%%%%%%%%%%
\section{En anglais}
%%%%%%%%%%%%%%%%%%%%%%%%%{\it }

Hendel, {\it Studies on the philosophy of D. Hume}, Princeton,
1925.

A. E. Taylor, {\it David Hume and the miraculous}, Gambridge, 1927.

J. Laird, {\it Hume’s Philosophy of Human Nature}, London,
1932.

{\it Hume and present day problems}, Aristotelian Society,
suppl., vol. XVIII, London, 1939 (4 symposia sur
l'identité du moi, sur les concepts {\it a priori}, sur l’éthique,
sur la religion naturelle avec des articles de Taylor,
de Lairp, de Jessop).

Norman Kemp Smith, {\it Philosophy of David Hume}, London,
1941.

D. G. C. Mac Nabb, {\it David Hume, His theory of Knowledge
and Morality}, London, 1954.

E. G. Mossner, {\it The life of David Hume}, London, 1954.

%%%%%%%%%%%%%%%%%%%%%%%%%
\section{En français}
%%%%%%%%%%%%%%%%%%%%%%%%%

G. Compayré, {\it }La philosophie de D. Hume, Paris, 1873.

G. Lechartier, {\it David Hume sociologue et moraliste},
Paris, 1900.

L. Levy-Bruhl, Orientation de la pensée de D. Hume,
{\it Revue de métaphysique et de morale}, 1909.

A. Leroy, {\it Critique et religion chez D. Hume}, Paris, 1931.

%92
Laporte, Le septicisme du Hume, {\it Revue philosophique,}
1933-1934.

G. Brercer, Husserl et Hume, {\it Revue internationale de
philosophie}, 1939.

{\it Mélange David Hume}, divers articles, {\it Revue internationale
de philosophie}, Bruxelles, 1952.

DELEUZE, Empirisme et subjectivité, Paris, Presses universitaires 
de France, 1953.

A. Leroy, {\it David Hume}, Paris, Presses Universitaires de
France, 1953.

O. Brunet, Philosophie et Esthétique chez D. Hume,
Paris, Librairie A.-G. Nizet, 1965.

%%%%%%%%%%%%%%%%%%%%%%%%%%%%%%%%%%%%%%%%%%%%%%%%%%%%%%%%%%%%%%%%%%%%%%%%

%
%
%%%%%%%%%%%%%%%%%%%%%
\section{Encyclopédie de la philosophie}
%%%%%%%%%%%%%%%%%%%%%
%
\subsection{Nécessité}
Est dit nécessaire « ce qui ne peut être autrement qu'il n'est »,
telle est du moins la définition aristotélicienne de la notion (grec {\it anankaion}).
Elle s'applique
notamment pour caractériser la nature de
la relation qui relie entre elles les propositions (prémisses et protases) du syllogisme scientifique ou démonstratif. Est
nécessaire ce qui appartient à un sujet
partout et toujours : ainsi la propriété
d’avoir la somme de ses angles intérieurs
égale à deux droits est une propriété
nécessaire pour tout triangle en tant que
tel. Mais, outre le sens logique et ontologique, il y a un sens psychologique, de
contrainte inévitable. Ce qui a conduit certains à émettre l'hypothèse que le concept
est peut-être né d’une projection anthropomorphique de l’idée originelle de coercition,
ce qui aurait pour conséquence que la
notion de nécessité déontologique (c’est-à-dire d’obligation) précéderait historiquement toutes les autres formes de nécessité.
La notion de nécessité physique et causale comme subordination aux lois de la
nature et celle de nécessité logique
comme propriété de ce qui est « forcément » vrai en vertu des lois logiques
seraient apparues par la suite, selon un
ordre inversé par rapport à celui qui est
parfois considéré logiquement rationnel.

\subsection{Nécessité et modalité}

La notion leibnizienne de nécessité
comme vérité dans tous les mondes possibles à été reprise dans l’analyse sémantique contemporaine des modalités.
D'après ces analyses, il apparaît que la
notion modale fondamentale est la notion
de possibilité plutôt que celle de nécessité. À côté de cette notion leibnizienne
de nécessité absolue, la logique. contemporaine a caractérisé différentes notions
de nécessité relative, en-isolant des sous-classes pertinentes à l’intérieur de la
classe de tous les mondes possibles,
comme celle des mondes où sont valides
les lois de la nature ou les lois de code
pénal. Toute logique de la nécessité,
quelle qu’elle soit, a comme condition
minimum d'éviter ce que l’on appelle
« l’effondrement des modalités », à savoir
la démonstration de l’équivalence entre
une proposition nécessaire et une proposition dépourvue d’opérateurs modaux
(cette équivalence est présente dans toute
philosophie fataliste ou strictement déterministe : il suffit de penser à l’assimilation
opérée par Hegel entre ce qui est rationnel et ce qui est réel). Dans une acception
élargie, la nécessité logique est parfois
assimilée à l’analycité. Les énoncés analytiques résultent nécessaires dans la
mesure où les règles linguistiques
« créent » la signification des termes que
l’on y trouve, que ces termes soient des
constantes logiques (« et », « ou »,
« tous », etc.) ou des mots du langage
ordinaire comme « célibataire » ou « non-
marié ». Du point de vue pragmatiste
d'auteurs comme W.V.O. Quine, la nécessité d’un énoncé quelconque consiste dans
son immunité à l’intérieur du système des
connaissances prouvées, c’est-à-dire dans
le fait que le renoncement à l’énoncé en
question est trop coûteux pour le système
dans son ensemble. Tout en excluant qu’il
existe une distinction historiquement définitive entre analytique et synthétique et
entre nécessaire et contingent, Quine.
pense toutefois que la nécessité logique
peut être justement attribuée. à des
énoncés dont la vérité dépend uniquement des constantes logiques qui sont
présentes en lui. Il s’agit d’un pas en avant
par rapport à la conception du Tractatus
de L. Wittgenstein, qui assimilait les
nécessités logiques aux tautologies du calcul propositionnel, excluant de cette
manière les vérités logiques qui dépendent de la présence de quantificateurs. De
toute façon, cette conviction est partagée
par presque toute  l'épistémologie
contemporaine, surtout par l’épistémologie empiriste : que la nécessité des lois
scientifiques et des inférences garanties
par elles ne dépend pas de l’existence de
connexions nécessaires dans la nature,
mais doit être indirectement référée à la
nécessité logique ou au concept d’implication logiquement nécessaire.

%\vspace{0.35cm}
%$\to$ existence ; mondes possibles ; possibilité ; quantificateurs ; Quine

%%%%%%%%%%%%%%%%%%%%%%%%%%%%%%%%%%%%%%%%%%%%%%%%%%%%%%%%%%%%%%%%%%%%%%%%

%%%%%%%%%%%%%%%%%%%%%%%%%%%%%%%%%%%%%
\section{Pratique de la philosophie}
%%%%%%%%%%%%%%%%%%%%%%%%%%%%%%%%%%%%

\subsection{Opinion}

{\footnotesize
\begin{itemize}[leftmargin=1cm, label=\ding{32}, itemsep=1pt]
\item {\bf \textsc{Étymologie} :} latin {\it opinari},
« émettre une opinion ».
\item {\bf \textsc{Sens ordinaire} :} avis,
jugement porté sur
un sujet, qui ne relève pas d'une
connaissance rationnelle vérifiable,
et dépend donc du système de
valeurs en fonction duquel on se
prononce.
\item {\bf \textsc{Philosophie} :} jugement
sans fondement rigoureux,
souvent dénoncé dans la mesure où
il se donne de façon abusive les
apparences d’un savoir.
\end{itemize}
}

L'interrogation sur la nature de la vérité
et les moyens de l’atteindre a conduit
nombre de philosophes à distinguer,
entre les différents types de connaissance
possibles, ceux qui conduisent effectivement
à la vérité, et ceux qui en éloignent.
En un premier sens, l’opinion est ainsi
traditionnellement considérée comme un
genre de connaissance peu fiable, fondée
sur des impressions, des sentiments, des
croyances où des jugements de valeur
subjectifs. Pour Spinoza, par exemple,
elle est forcément « sujette à l'erreur et n’a
jamais lieu à l'égard de quelque chose
dont nous sommes certains mais à l'égard
de ce que l’on dit conjecturer ou supposer »
({\it Court traité}, chap. II). Depuis Platon,
et jusque chez de nombreux penseurs
contemporains, l'opinion est
dénoncée comme a priori douteuse, illusoire
ou fausse, voire dangereuse, lorsqu’elle
cherche à s'imposer en dissimulant
la faiblesse de ses fondements sous
les apparences de la plus claire certitude.
Selon Adorno ({\it Modèles critiques}, 1963),
« l'opinion s’approprie ce que la connaissance
ne peut atteindre pour s’y substituer »,
elle rassure à bon compte, parce
qu’« elle offre des explications grâce auxquelles
on peut organiser sans contradiction
la réalité contradictoire ». Tel est bien
le « fonctionnement psychique » qui soustend,
par exemple, les opinions racistes :
pour être plus crédible, la peur de l’autre
prend le masque de l'affirmation de son
infériorité ou de la mise en garde contre
le danger qu'il est censé représenter. La
justesse de ces analyses ne doit pas faire
oublier qu'en un autre sens, l'opinion
constitue une forme de connaissance
utile, voire un type de jugements éminemment
respectables. Dans le {\it Ménon},
Platon reconnaît aux opinions droites la
faculté, sur les sujets qui ne relèvent ni de
la science ni de la simple conjecture,
d'éclairer l’action humaine. Dans le
domaine moral par exemple, à défaut de
vérités certaines, des intuitions justes
relatives au bien peuvent guider efficacement
l'éducation ou l’action, en leur
fixant pour but la satisfaction d'intérêts
conformes aux exigences de la réflexion,
et non à la soumission aux apparences ou
au plaisir immédiat. Enfin, sur toutes les
questions qui engagent des choix individuels
qu'aucune autorité ne peut légitimement
contraindre {\bf --} la religion, la
préférence politique, l'adhésion à une
conception du monde {\bf --} la liberté d’opinion
est un droit fondamental, dans les
sociétés démocratiques en tout cas, dès
l'instant où ceux auxquels elle est garantie
n'en usent pas au détriment de la
liberté d'autrui.

Analysée dans le {\it Traité
théologico-politique}, où Spinoza insiste
sur la nécessité d'une indépendance
absolue des opinions religieuses et de
leur expression par rapport à l'État, la
liberté d'opinion est proclamée dans la
Déclaration des droits de l'homme et du
citoyen de 1789. Et depuis près d'un
siècle, elle est au cœur du principe de la
laïcité qui garantit (en particulier en
France) la séparation entre l'Église et
l'État.

{\footnotesize
\begin{itemize}[leftmargin=1cm, label=\ding{32}, itemsep=1pt]
\item {\bf \textsc{Termes voisins} :} avis ; croyance.
\item {\bf \textsc{Termes opposés} :} science.
\end{itemize}
}

\subsubsection{Opinion publique}

Ensemble fluctuant de prises de positions
portant sur des questions politiques,
 morales, économiques... Les
« sondages d'opinion » prétendent en
constituer une sorte de baromètre.

{\footnotesize
\begin{itemize}[leftmargin=1cm, label=\ding{32}, itemsep=1pt]
\item {\bf \textsc{Corrélats} :} connaissance ;
conviction ; croyance ; doute ; foi ;
jugement ; préjugé.
\end{itemize}
}

%%%%%%%%%%%%%%%%%%%%%%%%%%%%%%%%%%%%%%%%%%%%%%%%%%%%
\subsection{Préjugé}

{\footnotesize
\begin{itemize}[leftmargin=1cm, label=\ding{32}, itemsep=1pt]
\item {\bf \textsc{Étymologie} :} latin {\it praejudicare},
« juger préalablement ».
\item {\bf \textsc{Sens ordinaire} :} Opinion admise sans
jugement ni raisonnement.
\end{itemize}
}

Le terme préjugé est souvent employé
dans un sens péjoratif, pour dénoncer
l'erreur ou au moins l'absence de
réflexion qui conduit un individu à
adhérer à une idée fausse {\bf --} dont il n’a
pas pris la peine de contrôler le bien-fondé {\bf --}
voire à la défendre contre des
idées justes, ou à condamner des individus
au nom de cette idée (par
exemple, les opinions racistes sont des
préjugés).

{\footnotesize
\begin{itemize}[leftmargin=1cm, label=\ding{32}, itemsep=1pt]
\item {\bf \textsc{Termes voisins} :} opinion.
\item {\bf \textsc{Termes opposés} :} savoir ; science.
\item {\bf \textsc{Corrélats} :} certitude ; croyance ;
dogme ; doute ; foi.
\end{itemize}
}

%%%%%%%%%%%%%%%%%%%%%%%%%%%%%%%%%%%%%%%%%%%%%%%%%%%%

\subsection{Erreur}

{\footnotesize
\begin{itemize}[leftmargin=1cm, label=\ding{32}, itemsep=1pt]
\item {\bf \textsc{Étymologie} :} latin {\it error}, « course
à l'aventure », de {\it errare}, «errer».
\item {\bf \textsc{Logique et sciences} :} affirmation
fausse, c'est-à-dire non conforme
aux règles de la logique, et/ou en
contradiction avec les données
expérimentales.
\item {\bf \textsc{Psychologie} :} état
de l'esprit qui tient pour vrai ce
qui est faux, et réciproquement
(ex. : « être dans l’erreur »).
\end{itemize}
}

L'erreur doit être soigneusement distinguée
aussi bien de la faute (qui engage
plus nettement notre responsabilité) que
de l’illusion (qui n’est pas vaincue par
le savoir). L'erreur procède toujours de
notre jugement : elle résulte, selon Descartes,
d’un décalage permanent entre
notre volonté, qui est infinie, et notre
entendement, qui ne l'est pas. Nous
nous trompons parce que nous outrepassons
nos possibilités intellectuelles,
par étourderie ou vanité : l'erreur n’est
donc qu'une privation de connaissance.
L'épistémologie contemporaine, au
contraire, donne à l’erreur un tout autre
statut, plus « positif ». Bachelard, notamment,
montre que les « vérités » scientifiques
ne sont jamais que provisoires,
qu'elles doivent constamment être remaniées
et corrigées. La connaissance
scientifique ne peut pas faire l'économie
de l’erreur.

{\footnotesize
\begin{itemize}[leftmargin=1cm, label=\ding{32}, itemsep=1pt]
\item {\bf \textsc{Termes voisins} :} fausseté ; illusion ;
incorrection.
\item {\bf \textsc{Termes opposés} :} vérité.
\item {\bf \textsc{Corrélats} :} connaissance ; Évidence ;
faute ; illusion ; jugement.
\end{itemize}
}

%%%%%%%%%%%%%%%%%%%%%%%%%%%%%%%%%%%%%%%%%%%%%%%%%%%%

%
%\newpage
%
\end{appendix}
%

%====================== INCLUSION DE LA BIBLIOGRAPHIE ======================
%
%récupérer les citation avec "/footnotemark" : 
%%\nocite{*}
%
% choix du style de la biblio
%%\bibliographystyle{plain}
%
% inclusion de la biblio
%%\cleardoublepage
%%\addcontentsline{toc}{chapter}{Bibliographie}
%%\bibliography{bibliographie.bib}
%
%====================== FIN DU DOCUMENT ======================
%
\end{document}
%%%%%%%%%%%%%%%%%%%%%%%%%%%%%%%%%%%%%%%%%%%%%%%%%%%%%%%%%%%%%%%%%%%%%%%%%%%%%%%%%
